\chapter{函数、极限与连续}\label{ch:1}
\begin{tcolorbox}[title=章首知识地图,colframe=Ink]
函数与定义域 $\to$ 极限语言 $\to$ 运算法则与无穷小 $\to$ 数列收敛 $\to$ 连续性及存在性。计算前先分清数列、函数、单侧与双侧极限。
\end{tcolorbox}
\begin{examtip}考研主线：极限计算、分段连续、递推数列和参数抵消。常见题型：先判断存在性，再确定参数或计算极限。竞赛拓展：Stolz、积分型极限和渐近误差；先证明收敛再求值。上述为学习任务分层，具体范围以所报考科目的当年大纲为准。\end{examtip}
\section{函数与极限的语言}\label{sec:1:1}

\begin{center}
\begin{tikzpicture}[scale=.65]
\draw[->,Ink] (-3.6,0)--(3.7,0) node[right]{$x$};\draw[->,Ink] (0,-.4)--(0,2.1) node[above]{$y$};
\draw[Teal,thick,domain=-3.2:3.2,samples=80] plot (\x,{1+0.12*\x*\x});
\fill[Teal!10] (-1.2,.8) rectangle (1.2,1.2);\draw[Teal,dashed](-1.2,.8)--(1.2,.8);\draw[Teal,dashed](-1.2,1.2)--(1.2,1.2);
\draw[Teal,thick,domain=-1.2:1.2,samples=40] plot (\x,{1+0.12*\x*\x});
\node[left] at (0,1.5){$L$附近};\node[below]at(1.2,0){$\delta$};\node[below]at(-1.2,0){$-\delta$};
\end{tikzpicture}\par\small 极限的误差控制：缩小水平邻域，使函数值进入指定误差带。
\end{center}

\subsection{概念、定理与推导}

单侧极限只允许从定义点的一边趋近；双侧极限存在当且仅当两侧有限极限存在且相等。$x\to\infty$时以$x>M$代替邻域条件；$f(x)\to+\infty$要求对任意$K>0$最终有$f(x)>K$。无穷大是变化状态，不是可参与普通四则运算的实数。单调性、奇偶性、周期性及反函数定义域应先于计算检查。


\begin{definition}{函数、数列与函数极限}函数由定义域、对应关系和陪域给定；复合$f\circ g$要求$g(x)$落在$f$的定义域。单射区分不同输入，双射才在整个陪域上有反函数。$a_n\to A$指任意$\varepsilon>0$，存在$N$使$n>N$时$|a_n-A|<\varepsilon$。设$a$是$f$的定义域$D$的聚点，即任意去心邻域均与$D$相交；$\lim_{x\to a}f(x)=L$指任意$\varepsilon>0$存在$\delta>0$，使$x\in D$且$0<|x-a|<\delta$蕴含$|f(x)-L|<\varepsilon$。聚点前提保证此定义不是对空邻域作判断，并保证极限的唯一性。\end{definition}
\begin{theorem}{极限性质与四则运算}有限极限唯一，收敛数列有界，函数有有限极限则在充分小的去心邻域内有界。若$f\to A,g\to B$，则$f\pm g\to A\pm B,fg\to AB$；当$B\ne0$时$\frac{f}{g}\to \frac{A}{B}$。若$f\le h\le g$且$f,g\to L$，则$h\to L$。\end{theorem}
\begin{proof}唯一性由$|L_1-L_2|\le|L_1-f|+|f-L_2|$与任意小的误差矛盾得到。乘法将$fg-AB$拆为$f(g-B)+B(f-A)$并使用局部有界性；夹逼直接比较$h-L$的上下界。\end{proof}
奇偶性要求定义域关于原点对称；周期$T$必须满足平移后仍在定义域且$f(x+T)=f(x)$。这些结构可在求极限前简化表达式。

\begin{method}{函数与极限的语言}条件：先确认定义域和趋近方向。步骤：检查直接代入；出现未定式再变形；用两个趋近序列可否定极限。易错点：函数在点上无定义不妨碍该点极限存在。\end{method}
\subsection{例题与方法比较}
\begin{example}{基础例题｜定义证明}{难度：$\star$}
\label{ex:1:1}\textbf{题目} 用定义证明$\lim_{x\to2}(3x-1)=5$。\par
\textbf{思路} 把函数误差换成自变量误差。\par
\textbf{解答} $|(3x-1)-5|=3|x-2|$。给定$\varepsilon>0$，取$\delta=\frac{\varepsilon}{3}$，即可使误差小于$\varepsilon$。另一方法：连续函数的四则法则给出$3\cdot2-1=5$。\par
\textbf{方法点评} 定义法提供误差控制；四则法适于计算。\par
\textbf{易错点或变式} 定义中的$\delta$可依赖$\varepsilon$，不能依赖正在变化的$x$。
\end{example}
\begin{example}{基础例题｜有理化}{难度：$\star$}
\label{ex:1:2}\textbf{题目} 求$\lim_{x\to0}\frac{\sqrt{1+x}-1}{x}$。\par
\textbf{思路} 消除相消或识别导数。\par
\textbf{解答} 方法一：乘共轭后变为$\frac{1}{\sqrt{1+x}+1}\to\frac{1}{2}$。方法二：这是$f(x)=\sqrt{1+x}$在0处的导数，$f\prime(0)=\frac{1}{2}$。\par
\textbf{方法点评} 有理化不需要预先掌握求导，导数定义便于识别结构。\par
\textbf{易错点或变式} 只在$x\ne0$时约去$x$，这与去心极限一致。
\end{example}
\begin{example}{考研风格例题｜不存在}{难度：$\star\star$}
\label{ex:1:3}\textbf{题目} 判断$\lim_{x\to0}\sin(\frac{1}{x})$。\par
\textbf{思路} 找两个极限值不同的趋近序列。\par
\textbf{解答} 取$x_n=(2n\pi+\frac{\pi}{2})^{-1}$和$y_n=(2n\pi+\frac{3\pi}{2})^{-1}$，两者趋零但函数值分别为1与$-1$，故极限不存在。另一思路：任意去心邻域的函数振幅均为2，不满足极限的Cauchy条件。\par
\textbf{方法点评} 序列判别给出可检验反例，振幅语言揭示振荡原因。\par
\textbf{易错点或变式} 有界不等于收敛。
\end{example}
\section{无穷小比较与重要极限}\label{sec:1:2}

\begin{theorem}{Heine序列判别，标准结论}设$a$是定义域的聚点。$\lim_{x\to a}f(x)=L$等价于：对任意定义域内的数列$x_n\ne a$且$x_n\to a$，均有$f(x_n)\to L$。\end{theorem}
\begin{proof}必要性把极限定义的$\delta$作用于收敛数列即可。若函数极限不是$L$，存在$\varepsilon_0>0$，使每个$\frac{1}{n}$邻域都能选出$x_n\ne a$且$|f(x_n)-L|\ge\varepsilon_0$，这与序列条件矛盾。\end{proof}
\begin{method}{先审查定义域，再用初等函数结构}
对数要求真数正，分母要求非零，偶次根要求根下非负。复合函数的定义域是所有同时满足内外层条件的点。反函数交换定义域与值域，例如$\arcsin$的定义域为$[-1,1]$，值域为$[-\frac{\pi}{2},\frac{\pi}{2}]$；反正切取值$(-\frac{\pi}{2},\frac{\pi}{2})$。偶函数双侧极限可由一侧推知；奇函数在0有有限极限时该极限只能为0，但奇性本身不保证极限存在。
\end{method}

\subsection{概念、定理与推导}

\begin{definition}{无穷小阶}在同一极限过程中，$f=o(g)$表示$\frac{f}{g}\to0$；$f=O(g)$表示比值有界；$f\sim g$表示$\frac{f}{g}\to1$，要求分母在充分小去心邻域不为零。\end{definition}
\begin{theorem}{两个重要极限}以弧度计，$\frac{\sin x}{x}\to1$；$\lim_{x\to0}(1+x)^{\frac{1}{x}}=\e$。因此$\sin x\sim x$、$\tan x\sim x$、$1-\cos x\sim \frac{x^2}{2}$、$\ln(1+x)\sim x$、$\e^x-1\sim x$。\end{theorem}
\begin{proof}在$0<x<\frac{\pi}{2}$，几何比较给$\sin x<x<\tan x$，夹逼后用偶性扩展到左侧。定义$\ln t=\int_1^t\frac{\dd u}{u}$，积分估计给$\frac{\ln(1+x)}{x}\to1$，取指数即得第二极限。\end{proof}
\begin{corollary}{等价替换的边界}乘积或商中可逐因子替换等价无穷小。和差中必须比较总误差：若$u-U=o(h),v-V=o(h)$，才有$u-v=(U-V)+o(h)$。\end{corollary}

\begin{method}{无穷小比较与重要极限}条件：各无穷小沿同一过程趋零。步骤：先提取主因子，再判定相消阶；指数型先取对数。易错点：$\sin x-x$不能用$x-x$代替。\end{method}
\subsection{例题与方法比较}
\begin{example}{考研风格例题｜等价替换}{难度：$\star\star$}
\label{ex:1:4}\textbf{题目} 求$\lim_{x\to0}\frac{\ln(1+2x)}{\sin3x}$。\par
\textbf{思路} 分子分母各有一阶主项。\par
\textbf{解答} 方法一：等价替换得$\frac{2x}{3x}\to\frac{2}{3}$。方法二：洛必达给$\frac{\frac{2}{1+2x}}{3\cos3x}\to\frac{2}{3}$。\par
\textbf{方法点评} 等价法短；洛必达需确认$\frac{0}{0}$与分母导数非零。\par
\textbf{易错点或变式} 定义域要求$1+2x>0$，零点附近满足。
\end{example}
\begin{example}{考研风格例题｜指数未定式}{难度：$\star\star$}
\label{ex:1:5}\textbf{题目} 求$\lim_{x\to0}(1+2x)^{\frac{3}{x}}$。\par
\textbf{思路} 先研究对数。\par
\textbf{解答} 方法一：对数为$\frac{3\ln(1+2x)}{x}\to6$，结果为$\e^6$。方法二：令$t=2x$，原式等于$[(1+t)^{\frac{1}{t}}]^6\to\e^6$。\par
\textbf{方法点评} 取对数处理变化指数，标准极限适合线性内层。\par
\textbf{易错点或变式} 实数幂要求底数在邻域内为正。
\end{example}
\begin{example}{基础例题｜夹逼}{难度：$\star$}
\label{ex:1:6}\textbf{题目} 求$\lim_{x\to0}x^2\sin(\frac{1}{x})$。\par
\textbf{思路} 把振荡因子视为有界量。\par
\textbf{解答} $|x^2\sin(\frac{1}{x})|\le x^2\to0$，故极限为0。它在0补值0后连续，但商$\sin(\frac{1}{x})$仍没有极限。\par
\textbf{方法点评} 有界量乘无穷小比机械求导更直接。\par
\textbf{易错点或变式} 不得将乘积极限拆成两个极限，因为第二个因子的极限不存在。
\end{example}
\section{洛必达与泰勒的选择}\label{sec:1:3}
\subsection{概念、定理与推导}

\begin{theorem}{洛必达法则，标准结论}在趋近端点的去心单侧邻域，$f,g$可导且$g\prime\ne0$；若$f,g\to0$，或$|g|\to\infty$，且$\frac{f\prime}{g\prime}\to L$（可为无穷），则$\frac{f}{g}\to L$。双侧结论须两侧一致。\end{theorem}
\begin{theorem}{Taylor展开}若$f$在$a$附近具有$n+1$阶连续导数，则
\begin{equation}f(a+h)=\sum_{k=0}^{n}\dfrac{f^{(k)}(a)}{k!}h^k+\dfrac{f^{(n+1)}(a+\theta h)}{(n+1)!}h^{n+1},\quad0<\theta<1.\end{equation}
特别地，$\e^x=1+x+\frac{x^2}{2}+\frac{x^3}{6}+O(x^4)$，$\sin x=x-\frac{x^3}{6}+O(x^5)$，$\ln(1+x)=x-\frac{x^2}{2}+\frac{x^3}{3}+O(x^4)$。\end{theorem}
\begin{proof}减去Taylor多项式，再减去常数倍$(t-a)^{n+1}$使函数在$a,x$处为零。多次应用Rolle定理得到余项。洛必达的$\frac{0}{0}$情形来自Cauchy中值定理；一般端点情形需先在截断区间应用，再取极限。\end{proof}

\begin{method}{洛必达与泰勒的选择}条件：洛必达检查可导与未定式，Taylor检查展开阶数。步骤：相消至哪一阶就展开至哪一阶；把余项一并除以分母。易错点：导数比值无极限不能推出原比值无极限。\end{method}
\subsection{例题与方法比较}
\begin{example}{考研风格例题｜三阶相消}{难度：$\star\star$}
\label{ex:1:7}\textbf{题目} 求$\lim_{x\to0}\frac{\sin x-x}{x^3}$。\par
\textbf{思路} 一阶项相消，保留三阶。\par
\textbf{解答} Taylor给$\frac{-\frac{x^3}{6}+O(x^5)}{x^3}\to-\frac{1}{6}$。洛必达连续用三次，依次得到$\frac{\cos x-1}{3x^2}$、$-\frac{\sin x}{6x}$、$-\frac{\cos x}{6}\to-\frac{1}{6}$。\par
\textbf{方法点评} Taylor揭示阶；洛必达可核验符号。\par
\textbf{易错点或变式} 每次求导前仍需确认当前式子是未定式。
\end{example}
\begin{example}{考研风格例题｜参数抵消}{难度：$\star\star$}
\label{ex:1:8}\textbf{题目} 确定$a,b$使$\lim_{x\to0}\frac{\e^x-a-bx}{x^2}$有限，并求极限。\par
\textbf{思路} 常数项、一阶项必须依次为零。\par
\textbf{解答} 展开分子为$(1-a)+(1-b)x+\frac{x^2}{2}+O(x^3)$。有限极限迫使$a=b=1$，此时极限$\frac{1}{2}$。若$a\ne1$则二阶分母不能平衡常数；若$a=1,b\ne1$则仍有$\frac{1}{x}$项。\par
\textbf{方法点评} 先分类再计算，比含参连续求导可靠。\par
\textbf{易错点或变式} 有限双侧极限比单侧无穷极限要求更强。
\end{example}
\begin{example}{考研风格例题｜差式}{难度：$\star\star$}
\label{ex:1:9}\textbf{题目} 求$\lim_{x\to0}(\frac{1}{x}-\frac{1}{\e^x-1})$。\par
\textbf{思路} 先合并成一个商。\par
\textbf{解答} 合并得$\frac{\e^x-1-x}{x(\e^x-1)}$。分子$\sim \frac{x^2}{2}$，分母$\sim x^2$，极限为$\frac{1}{2}$。\par
\textbf{方法点评} 合并把两个无穷量之差转成可比较的无穷小。\par
\textbf{易错点或变式} 分别求两个无穷量的极限再相减没有意义。
\end{example}
\section{数列、单调有界与Stolz}\label{sec:1:4}
\subsection{概念、定理与推导}

\begin{theorem}{单调有界与Cauchy准则}实数单调递增且有上界的数列收敛于其上确界。实数列收敛当且仅当任意$\varepsilon>0$存在$N$使$m,n>N$时$|a_m-a_n|<\varepsilon$。\end{theorem}
\begin{proof}设$A=\sup\{a_n\}$。存在$a_N>A-\varepsilon$，单调性给$A-\varepsilon<a_n\le A$。Cauchy准则的充分性依赖实数完备性，可由有界子列收敛再控制全列证明。\end{proof}
\begin{theorem}{Stolz定理的无穷分母型}若$b_n$严格递增且$b_n\to+\infty$，并且$\frac{a_{n+1}-a_n}{b_{n+1}-b_n}\to L$，则$\frac{a_n}{b_n}\to L$。\end{theorem}
\begin{proof}对充分大的$n$，把差分比夹在$L\pm\varepsilon$之间；乘以正差分后从固定$N$累加到$n-1$。除以$b_n$，固定初始项消失。无穷$L$同理。\end{proof}

\begin{method}{数列、单调有界与Stolz}条件：递推列先找不变区间，求和比先查分母严格单调无界。步骤：证明存在，再令极限满足递推方程。易错点：只解不动点方程并不能证明收敛。\end{method}
\subsection{例题与方法比较}
\begin{example}{基础例题｜递推}{难度：$\star$}
\label{ex:1:10}\textbf{题目} $a_1=1,a_{n+1}=\sqrt{2+a_n}$，求极限。\par
\textbf{思路} 证明$[1,2]$不变且递增。\par
\textbf{解答} 若$1\le a_n<2$，则$1<a_{n+1}<2$；又$2+a_n-a_n^2=(2-a_n)(a_n+1)>0$，故$a_{n+1}>a_n$。收敛极限满足$L^2=2+L$且$L\ge1$，故$L=2$。\par
\textbf{方法点评} 完整链条是有界、单调、连续传极限、筛根。\par
\textbf{易错点或变式} 舍去负根$-1$依据原数列非负。
\end{example}
\begin{example}{考研风格例题｜求和比}{难度：$\star\star$}
\label{ex:1:11}\textbf{题目} 求$\lim_{n\to\infty}\sum_{k=1}^n \frac{k^2}{n^3}$。\par
\textbf{思路} 可用差分、求和公式或Riemann和。\par
\textbf{解答} 方法一：差分比$\frac{(n+1)^2}{(n+1)^3-n^3}\to\frac{1}{3}$。方法二：平方和$\frac{n(n+1)(2n+1)}{6}$除以$n^3$。方法三：写成$\frac1n\sum(\frac{k}{n})^2\to\int_0^1t^2\dd t=\frac{1}{3}$。\par
\textbf{方法点评} 求和公式最快；Stolz不需闭式；积分法可推广连续函数。\par
\textbf{易错点或变式} 差分分母是$3n^2+3n+1$，不是$3n^2$的精确等式。
\end{example}
\begin{example}{考研风格例题｜根式平均}{难度：$\star\star$}
\label{ex:1:12}\textbf{题目} 求$\lim_{n\to\infty}\frac{(n!)^{\frac{1}{n}}}{n}$。\par
\textbf{思路} 对数化为Riemann和，但零点有对数奇性。\par
\textbf{解答} 令$L_n=\frac1n\sum_{k=1}^n\ln(\frac{k}{n})$。由$\ln t$递增，用区间积分上下夹逼得$\int_0^1\ln t\dd t\le L_n\le\int_{\frac{1}{n}}^1\ln t\dd t$；两边趋于$-1$。故所求为$\e^{-1}$。\par
\textbf{方法点评} 有奇点的Riemann和用单调积分估计补严谨性。\par
\textbf{易错点或变式} 不能直接援用闭区间连续函数的Riemann和定理。
\end{example}
\section{连续性与综合极限}\label{sec:1:5}
\subsection{概念、定理与推导}

\begin{definition}{连续、间断与一致连续}设$a\in D$。$f$在$a$连续指：任意$\varepsilon>0$都存在$\delta>0$，使$x\in D$且$|x-a|<\delta$时$|f(x)-f(a)|<\varepsilon$。当$a$为$D$的聚点时，这等价于相对定义域的$\lim_{x\to a}f(x)=f(a)$；区间端点只检查可趋近的一侧，孤立点自动连续。对区间内点的间断，左右极限有限但不等为跳跃间断；相等但点值缺失或与之不符为可去间断；至少一侧无有限极限为第二类间断。$f$在$D$上一致连续指：任意$\varepsilon>0$存在同一$\delta>0$，使任意$x,y\in D$只要$|x-y|<\delta$就有$|f(x)-f(y)|<\varepsilon$。\end{definition}
\begin{theorem}{闭区间连续函数}连续函数在$[a,b]$上有界，取得最大最小值，具有介值性且一致连续。若$f(a)f(b)<0$，则$(a,b)$内有零点。\end{theorem}
\begin{proof}最值与一致连续性可由反证构造点列，利用闭有界区间的子列收敛性及连续性导出矛盾。介值性由二分区间和完备性得到。\end{proof}
\begin{extension}积分型极限先把求和写成“步长乘函数值”；递推型极限先寻找不变区间。存在性证明与数值计算应分开完成。\end{extension}

\begin{method}{连续性与综合极限}条件：闭区间与连续缺一不可。步骤：判间断先左右极限，证明存在先建连续函数并找端点异号。易错点：开区间连续不保证最值。\end{method}
\subsection{例题与方法比较}
\begin{example}{竞赛风格例题｜积分极限}{难度：$\star\star\star$}
\label{ex:1:13}\textbf{题目} 求$\lim_{n\to\infty}\sum_{k=1}^n \frac{n}{n^2+k^2}$。\par
\textbf{思路} 抽出步长$\frac{1}{n}$。\par
\textbf{解答} 原和$=\frac1n\sum_{k=1}^n[1+(\frac{k}{n})^2]^{-1}\to\int_0^1\frac{\dd t}{1+t^2}=\frac{\pi}{4}$。用递减函数的左右矩形估计，误差不超过$\frac{1}{n}$。\par
\textbf{方法点评} 积分不仅求极限，还提供误差量级。\par
\textbf{易错点或变式} 不能把单项趋零当成总和趋零。
\end{example}
\begin{example}{竞赛风格例题｜不动点}{难度：$\star\star\star$}
\label{ex:1:14}\textbf{题目} 连续函数$f:[0,1]\to[0,1]$，证明存在$c$满足$f(c)=c$。\par
\textbf{思路} 将不动点改写为零点。\par
\textbf{解答} 令$g(x)=f(x)-x$，有$g(0)\ge0,g(1)\le0$。端点为零则完成；否则由介值定理得内部零点。\par
\textbf{方法点评} 构造差函数把映射条件转成符号条件。\par
\textbf{易错点或变式} 结论不保证唯一；$f(x)=x$的每点都是不动点。
\end{example}
\begin{example}{竞赛风格例题｜细化极限}{难度：$\star\star\star$}
\label{ex:1:15}\textbf{题目} 求$\lim_{n\to\infty}n[(1+\frac{1}{n})^n-\e]$。\par
\textbf{思路} 先展开对数，再回到指数。\par
\textbf{解答} $n\ln(1+\frac{1}{n})=1-\frac{1}{2n}+O(n^{-2})$，故$(1+\frac{1}{n})^n=\e[1-\frac{1}{2n}+O(n^{-2})]$。极限为$-\frac{\e}{2}$。\par
\textbf{方法点评} 先对数后指数能控制指数型误差。\par
\textbf{易错点或变式} 只知道底层极限为$\e$不足以确定一阶误差。
\end{example}
\begin{warning}四则法只在相应有限极限存在时使用；等价替换在差式中要控制误差；对递推式直接“令$n\to\infty$”必须先证明收敛。\end{warning}


\section{实数基础、映射与逼近}\label{mat:real_foundations}
\begin{examtip}完备性、映射与数值逼近构成极限方法的逻辑基础。\end{examtip}
\subsection{确界、区间套与完备性}\label{mat:real_foundations:complete}
\begin{method}{确界、区间套与完备性}\textbf{适用条件} 有界非空实数集或收缩闭区间。\par\textbf{实施与依据} 上确界是最小上界；区间套长度趋零时交集恰有一点。两者与实数Cauchy完备性等价。\end{method}
\begin{example}{专题例题｜确界、区间套与完备性}{难度：$\star\star$}
\textbf{题目} 用完备性构造$\sqrt2$。\par\textbf{思路与解答} 取$A=\{x\ge0:x^2<2\}$，非空有界，令$s=\sup A$。若$s^2<2$，可取充分小$h>0$使$(s+h)^2<2$，违背上界；若$s^2>2$，可取$h$使$(s-h)^2>2$，违背最小性。因此$s^2=2$。\par\textbf{方法点评与边界} Dedekind分割与Cauchy列等价类均可构造实数，本册采用确界公理组织证明。
\end{example}
\subsection{映射、反函数与函数方程}\label{mat:real_foundations:mapping}
\begin{method}{映射、反函数与函数方程}\textbf{适用条件} 复合、单射和对称条件。\par\textbf{实施与依据} 先核对定义域与值域，再用代入、对称与单射消元。\end{method}
\begin{example}{专题例题｜映射、反函数与函数方程}{难度：$\star\star$}
\textbf{题目} 连续函数$f:\R\to\R$满足$f(x+y)=f(x)+f(y)$，证明$f(x)=cx$。\par\textbf{思路与解答} 令$c=f(1)$。代0得$f(0)=0$，归纳得$f(n)=nc$，再由$f(\frac{m}{n})=\frac{mc}{n}$得有理点成立。对任意实数取有理数列趋近，用连续性得到$f(x)=cx$。\par\textbf{方法点评与边界} 若无连续或其他正则性条件，不能把有理点结论直接延伸到实数。
\end{example}
\subsection{有理逼近与无理测度}\label{mat:real_foundations:approx}
\begin{method}{有理逼近与无理测度}\textbf{适用条件} 需要用有理数逼近实数。\par\textbf{实施与依据} 将$0,\alpha,\ldots,N\alpha$的小数部分放入$N$等分区间，使用抽屉原理。\end{method}
\begin{example}{专题例题｜有理逼近与无理测度}{难度：$\star\star$}
\textbf{题目} 证明任意实数$\alpha$与整数$N\ge1$，存在$1\le q\le N$和整数$p$，使$|q\alpha-p|\le\frac{1}{N}$。\par\textbf{思路与解答} 两小数部分距离不超过$\frac{1}{N}$，相减得$q\alpha-p$，其中$q$为两指标差。无理数的无理测度可用满足$|\alpha-\frac{p}{q}|<q^{-\mu}$的无穷多逼近指数的上确界描述。\par\textbf{方法点评与边界} 抽屉论证给Dirichlet型逼近；不能由此声称任意无理数的最佳指数都相同。
\end{example}
\subsection{一致连续与反例}\label{mat:real_foundations:uniform}
\begin{method}{一致连续与反例}\textbf{适用条件} 全域统一误差控制。\par\textbf{实施与依据} 紧集连续推出一致连续；非紧集上需另查导数界或显式反例。\end{method}
\begin{example}{专题例题｜一致连续与反例}{难度：$\star\star$}
\textbf{题目} 证明$x^2$在$\R$不一致连续。\par\textbf{思路与解答} 取$x_n=n,y_n=n+\frac{1}{n}$，距离趋0但$y_n^2-x_n^2=2+\frac{1}{n^2}\to2$，违背统一连续性。它在每个闭有界区间上一致连续。\par\textbf{方法点评与边界} 连续是逐点条件，一致连续是全域条件。
\end{example}

\section{35种极限求法：完整方法链}\label{mat:limits_35}
\begin{examtip}以下按35类方法组织，每类给出适用条件、执行步骤和独立算例；相似方法仍保留不同入口。\end{examtip}
\subsection{第1法：定义证明}\label{mat:limits_35:k01}
\begin{method}{第1法：定义证明}\textbf{适用条件} 需要验证收敛或给出误差阈值。\par\textbf{实施与依据} 由$|a_n-A|<\varepsilon$反解$N$；函数极限反解$\delta$。\end{method}
\begin{example}{专题例题｜第1法：定义证明}{难度：$\star\star$}
\textbf{题目} 证明$\frac{2n+1}{n+3}\to2$。\par\textbf{思路与解答} 误差为$\frac{5}{n+3}$。取$N>\frac{5}{\varepsilon}$即可。定义证明还直接给误差上界。\par\textbf{方法点评与边界} 阈值不能依赖被考察的$n$。
\end{example}
\subsection{第2法：直接代入}\label{mat:limits_35:k02}
\begin{method}{第2法：直接代入}\textbf{适用条件} 函数在趋近点连续，表达式有意义。\par\textbf{实施与依据} 先验分母与对数根式条件，再用连续函数复合。\end{method}
\begin{example}{专题例题｜第2法：直接代入}{难度：$\star\star$}
\textbf{题目} 求$\lim_{x\to1}\frac{\ln(1+x^2)}{1+x}$。\par\textbf{思路与解答} 分母在1处为2，真数为2，连续复合给$\frac{\ln2}{2}$。\par\textbf{方法点评与边界} 出现$\frac{0}{0}$时已不满足直接代入的条件。
\end{example}
\subsection{第3法：单侧极限与分段函数}\label{mat:limits_35:k03}
\begin{method}{第3法：单侧极限与分段函数}\textbf{适用条件} 分段界点、绝对值或不同分支。\par\textbf{实施与依据} 分别求左右极限，并比较二者是否相等。\end{method}
\begin{example}{专题例题｜第3法：单侧极限与分段函数}{难度：$\star\star$}
\textbf{题目} 求使$f(x)=ax+1\ (x<0),\ \e^x+b\ (x\ge0)$在0连续的条件。\par\textbf{思路与解答} 左极限1，右极限及点值$1+b$，故$b=0$，$a$任意。连续性不要求左右斜率相等；可导还需$a=1$。\par\textbf{方法点评与边界} 点值只影响连续性，不改变去心极限。
\end{example}
\subsection{第4法：概念判断与增长阶}\label{mat:limits_35:k04}
\begin{method}{第4法：概念判断与增长阶}\textbf{适用条件} 有界乘无穷小或可直接比较增长。\par\textbf{实施与依据} 先将表达式分成有界因子与趋零因子，或取对数比较增长。\end{method}
\begin{example}{专题例题｜第4法：概念判断与增长阶}{难度：$\star\star$}
\textbf{题目} 求$\lim_{x\to\infty}x^3\e^{-x}$。\par\textbf{思路与解答} 对数$3\ln x-x\to-\infty$，故极限0；也可连续三次洛必达。\par\textbf{方法点评与边界} 指数压过固定幂，不能推广到指数参数也变化而未经比较的情形。
\end{example}
\subsection{第5法：子列判别}\label{mat:limits_35:k05}
\begin{method}{第5法：子列判别}\textbf{适用条件} 振荡数列或奇偶递推。\par\textbf{实施与依据} 两子列不同则发散；若奇偶两子列同趋$L$则全列趋$L$。\end{method}
\begin{example}{专题例题｜第5法：子列判别}{难度：$\star\star$}
\textbf{题目} 判断$a_n=(-1)^n+\frac{1}{n}$。\par\textbf{思路与解答} 偶子列趋1，奇子列趋$-1$，因此发散。作为变式，$\frac{(-1)^n}{n}\to0$。\par\textbf{方法点评与边界} 一个收敛子列不能说明全列收敛。
\end{example}
\subsection{第6法：Heine序列化}\label{mat:limits_35:k06}
\begin{method}{第6法：Heine序列化}\textbf{适用条件} 把函数极限转成任意趋近序列。\par\textbf{实施与依据} 证明时处理任意序列；反证时选特定坏序列。\end{method}
\begin{example}{专题例题｜第6法：Heine序列化}{难度：$\star\star$}
\textbf{题目} 证明$\lim_{x\to0}x\sin(\frac{1}{x})=0$。\par\textbf{思路与解答} 任取$x_n\ne0,x_n\to0$，有$|x_n\sin(\frac{1}{x_n})|\le|x_n|\to0$；Heine定理给函数极限。\par\textbf{方法点评与边界} 验证两条路径只足以否定，不足以肯定极限。
\end{example}
\subsection{第7法：因式分解}\label{mat:limits_35:k07}
\begin{method}{第7法：因式分解}\textbf{适用条件} 多项式商出现公共零因子。\par\textbf{实施与依据} 分解、在去心邻域约去公共因子，再代入。\end{method}
\begin{example}{专题例题｜第7法：因式分解}{难度：$\star\star$}
\textbf{题目} 求$\lim_{x\to1}\frac{x^3-1}{x^2-1}$。\par\textbf{思路与解答} 约去$x-1$得$\frac{x^2+x+1}{x+1}\to\frac{3}{2}$。洛必达也给$\frac{3x^2}{2x}\to\frac{3}{2}$。\par\textbf{方法点评与边界} 约去的因子只需在去心邻域非零。
\end{example}
\subsection{第8法：无穷大转无穷小}\label{mat:limits_35:k08}
\begin{method}{第8法：无穷大转无穷小}\textbf{适用条件} 自变量趋于无穷，适合倒数代换。\par\textbf{实施与依据} 令$t=\frac{1}{x}$，把过程转至$0^+$或$0$。\end{method}
\begin{example}{专题例题｜第8法：无穷大转无穷小}{难度：$\star\star$}
\textbf{题目} 求$\lim_{x\to\infty}x\ln(1+\frac{1}{x})$。\par\textbf{思路与解答} 令$t=\frac{1}{x}\to0^+$，变为$\frac{\ln(1+t)}{t}\to1$。\par\textbf{方法点评与边界} 代换后方向也要改变。
\end{example}
\subsection{第9法：有理化}\label{mat:limits_35:k09}
\begin{method}{第9法：有理化}\textbf{适用条件} 根式相减造成消去。\par\textbf{实施与依据} 乘共轭，必要时连续有理化两次。\end{method}
\begin{example}{专题例题｜第9法：有理化}{难度：$\star\star$}
\textbf{题目} 求$\lim_{n\to\infty}(\sqrt{n^2+n}-n)$。\par\textbf{思路与解答} 共轭后为$\frac{n}{\sqrt{n^2+n}+n}=\frac{1}{\sqrt{1+\frac{1}{n}}+1}\to\frac{1}{2}$。\par\textbf{方法点评与边界} 不是每个根式差都应展开；先试精确恒等变形。
\end{example}
\subsection{第10法：先求和、求积或恒等改写}\label{mat:limits_35:k10}
\begin{method}{第10法：先求和、求积或恒等改写}\textbf{适用条件} 有限和、积可裂项或望远镜。\par\textbf{实施与依据} 先获得精确部分和，再取极限。\end{method}
\begin{example}{专题例题｜第10法：先求和、求积或恒等改写}{难度：$\star\star$}
\textbf{题目} 求$\lim_{n\to\infty}\sum_{k=1}^n\frac{1}{k(k+1)}$。\par\textbf{思路与解答} 裂成$\frac{1}{k}-\frac{1}{k+1}$，和为$1-\frac{1}{n+1}$，极限1。\par\textbf{方法点评与边界} 不能把$n$项的单项极限相加代替整体极限。
\end{example}
\subsection{第11法：对数恒等式}\label{mat:limits_35:k11}
\begin{method}{第11法：对数恒等式}\textbf{适用条件} 变化指数、乘积或$1^\infty$。\par\textbf{实施与依据} 正底数下先研究对数；取指数前确认对数极限。\end{method}
\begin{example}{专题例题｜第11法：对数恒等式}{难度：$\star\star$}
\textbf{题目} 求$\lim_{x\to0}(\cos x)^{\frac{1}{x^2}}$。\par\textbf{思路与解答} 对数$\frac{\ln\cos x}{x^2}\to-\frac{1}{2}$，故结果$\e^{-\frac{1}{2}}$。其中$\ln\cos x\sim\cos x-1\sim-\frac{x^2}{2}$。\par\textbf{方法点评与边界} 零点附近$\cos x>0$，保证实数幂有定义。
\end{example}
\subsection{第12法：三角恒等变换}\label{mat:limits_35:k12}
\begin{method}{第12法：三角恒等变换}\textbf{适用条件} 和差、倍角或乘积可消去。\par\textbf{实施与依据} 优先和差化积与半角式，后用重要极限。\end{method}
\begin{example}{专题例题｜第12法：三角恒等变换}{难度：$\star\star$}
\textbf{题目} 求$\lim_{x\to0}\frac{1-\cos2x}{x\sin3x}$。\par\textbf{思路与解答} $1-\cos2x=2\sin^2x$，故极限为$\frac{2}{3}$。Taylor给同一主项比。\par\textbf{方法点评与边界} 三角函数角度须用弧度。
\end{example}
\subsection{第13法：两个重要极限}\label{mat:limits_35:k13}
\begin{method}{第13法：两个重要极限}\textbf{适用条件} 能化成$\frac{\sin t}{t}$或$(1+t)^{\frac{1}{t}}$。\par\textbf{实施与依据} 完整保留内层尺度和外部系数。\end{method}
\begin{example}{专题例题｜第13法：两个重要极限}{难度：$\star\star$}
\textbf{题目} 求$\lim_{n\to\infty}(1+\frac{3}{n})^{2n}$。\par\textbf{思路与解答} 令$t=\frac{3}{n}$，原式$[(1+t)^{\frac{1}{t}}]^6\to\e^6$。\par\textbf{方法点评与边界} 不能遗漏尺度3与指数2的乘积。
\end{example}
\subsection{第14法：一般变量替换}\label{mat:limits_35:k14}
\begin{method}{第14法：一般变量替换}\textbf{适用条件} 同一内层反复出现或趋近点不在0。\par\textbf{实施与依据} 用内层量为新变量，并确认可达方向。\end{method}
\begin{example}{专题例题｜第14法：一般变量替换}{难度：$\star\star$}
\textbf{题目} 求$\lim_{x\to\frac{\pi}{2}}\frac{1-\sin x}{(x-\frac{\pi}{2})^2}$。\par\textbf{思路与解答} 设$t=x-\frac{\pi}{2}$，$\sin x=\cos t$，结果$\frac{1}{2}$。\par\textbf{方法点评与边界} 函数替换只需对应趋近过程合法，不必全局一一。
\end{example}
\subsection{第15法：等价无穷小}\label{mat:limits_35:k15}
\begin{method}{第15法：等价无穷小}\textbf{适用条件} 乘除结构且总误差不被相消放大。\par\textbf{实施与依据} 逐个因子化成已知主项；差式先合并或展开。\end{method}
\begin{example}{专题例题｜第15法：等价无穷小}{难度：$\star\star$}
\textbf{题目} 求$\lim_{x\to0}\frac{\tan2x\,\ln(1+3x)}{x^2}$。\par\textbf{思路与解答} 两因子分别$\sim2x,3x$，极限6。\par\textbf{方法点评与边界} 不能把$\tan x-\sin x$逐项换成$x-x$。
\end{example}
\subsection{第16法：洛必达与Stolz差分}\label{mat:limits_35:k16}
\begin{method}{第16法：洛必达与Stolz差分}\textbf{适用条件} 导数法针对函数商，差分法针对递增无界分母的数列商。\par\textbf{实施与依据} 未定式先规范化，每次检查所用法则条件。\end{method}
\begin{example}{专题例题｜第16法：洛必达与Stolz差分}{难度：$\star\star$}
\textbf{题目} 求$\lim\frac{\ln n}{n}$与$\lim_{n\to\infty}\sum_{k=1}^n\frac{k}{n^2}$。\par\textbf{思路与解答} 第一式用函数洛必达得$\frac{1}{x}\to0$。第二式差分比$\frac{n+1}{2n+1}\to\frac{1}{2}$，由Stolz得到极限。\par\textbf{方法点评与边界} 不能对一般数列直接写导数；须先找到合法函数模型。
\end{example}
\subsection{第17法：夹逼准则}\label{mat:limits_35:k17}
\begin{method}{第17法：夹逼准则}\textbf{适用条件} 精确计算困难，但可建立同极限上下界。\par\textbf{实施与依据} 找统一界，控制整体误差。\end{method}
\begin{example}{专题例题｜第17法：夹逼准则}{难度：$\star\star$}
\textbf{题目} 求$\lim_{n\to\infty}(n^2+2^n)^{\frac{1}{n}}$。\par\textbf{思路与解答} 最终$n^2\le2^n$，所以$2\le(n^2+2^n)^{\frac{1}{n}}\le2\cdot2^{\frac{1}{n}}$，极限2。\par\textbf{方法点评与边界} 比较指数速度时应说明不等式从充分大$n$起成立。
\end{example}
\subsection{第18法：单调有界与压缩映射}\label{mat:limits_35:k18}
\begin{method}{第18法：单调有界与压缩映射}\textbf{适用条件} 递推迭代或不动点。\par\textbf{实施与依据} 找不变区间，证明单调或$|T(u)-T(v)|\le q|u-v|$，$q<1$。\end{method}
\begin{example}{专题例题｜第18法：单调有界与压缩映射}{难度：$\star\star$}
\textbf{题目} $a_1=0,a_{n+1}=\cos a_n$，证明收敛。\par\textbf{思路与解答} 从第二步起$a_n\in[0,1]$。该区间$|T\prime|\le\sin1<1$，故相邻差受几何级数控制，数列为Cauchy，趋于唯一方程$L=\cos L$的解。\par\textbf{方法点评与边界} 该列不必单调，压缩法比单调法更适合。
\end{example}
\subsection{第19法：微分、积分中值定理}\label{mat:limits_35:k19}
\begin{method}{第19法：微分、积分中值定理}\textbf{适用条件} 小区间函数差或平均值。\par\textbf{实施与依据} 把函数差写成导数乘长度；把积分写成函数中间值乘区间长。\end{method}
\begin{example}{专题例题｜第19法：微分、积分中值定理}{难度：$\star\star$}
\textbf{题目} 求$\lim_{x\to0}\frac1x\int_x^{2x}\e^{t^2}\dd t$。\par\textbf{思路与解答} 积分中值为$\e^{\xi_x^2}x$，$\xi_x$在$x,2x$之间，所以极限1。\par\textbf{方法点评与边界} $x<0$时带方向积分仍与区间长度符号一致。
\end{example}
\subsection{第20法：Taylor主项}\label{mat:limits_35:k20}
\begin{method}{第20法：Taylor主项}\textbf{适用条件} 多项低阶抵消。\par\textbf{实施与依据} 展开到首个非零项，余项除以分母后仍须趋零。\end{method}
\begin{example}{专题例题｜第20法：Taylor主项}{难度：$\star\star$}
\textbf{题目} 求$\lim_{x\to0}\frac{\ln(1+x)-x+\frac{x^2}{2}}{x^3}$。\par\textbf{思路与解答} 分子$\frac{x^3}{3}+O(x^4)$，极限$\frac{1}{3}$；三次洛必达可核对。\par\textbf{方法点评与边界} 展开多一阶可用于判断误差符号和收敛速度。
\end{example}
\subsection{第21法：导数定义}\label{mat:limits_35:k21}
\begin{method}{第21法：导数定义}\textbf{适用条件} 差商或可整理为差商。\par\textbf{实施与依据} 辨认中心点与增量，不必机械展开。\end{method}
\begin{example}{专题例题｜第21法：导数定义}{难度：$\star\star$}
\textbf{题目} 求$\lim_{h\to0}\frac{(1+h)^\alpha-1}{h}$。\par\textbf{思路与解答} 这是$t^\alpha$在1处导数，结果$\alpha$。对数展开也给$(1+h)^\alpha=1+\alpha h+o(h)$。\par\textbf{方法点评与边界} $\alpha$固定；若同时依赖$h$，需另作分析。
\end{example}
\subsection{第22法：定积分与重积分定义}\label{mat:limits_35:k22}
\begin{method}{第22法：定积分与重积分定义}\textbf{适用条件} 求和具有步长和采样点。\par\textbf{实施与依据} 一维抽出$\frac{1}{n}$，二维抽出$\frac{1}{n^2}$。\end{method}
\begin{example}{专题例题｜第22法：定积分与重积分定义}{难度：$\star\star$}
\textbf{题目} 求$\lim_{n\to\infty}n^{-2}\sum_{i,j=1}^n(\frac{i}{n}+\frac{j}{n})^2$。\par\textbf{思路与解答} 为$\int_0^1\int_0^1(x+y)^2\dd x\dd y=\frac{1}{3}+\frac{1}{2}+\frac{1}{3}=\frac{7}{6}$。\par\textbf{方法点评与边界} 有奇点时需另加端点估计，不能照搬连续函数版本。
\end{example}
\subsection{第23法：变上限积分极限}\label{mat:limits_35:k23}
\begin{method}{第23法：变上限积分极限}\textbf{适用条件} 分子为积分，端点随参数趋近。\par\textbf{实施与依据} 用基本定理求导，或局部缩放$t=xu$。\end{method}
\begin{example}{专题例题｜第23法：变上限积分极限}{难度：$\star\star$}
\textbf{题目} 求$\lim_{x\to0}x^{-3}\int_0^x\sin(t^2)\dd t$。\par\textbf{思路与解答} 洛必达得$\frac{\sin(x^2)}{3x^2}\to\frac{1}{3}$。缩放后为$\int_0^1[\frac{\sin(x^2u^2)}{x^2}]\dd u\to\int_0^1u^2\dd u$。\par\textbf{方法点评与边界} 第二法需统一界$|\frac{\sin(x^2u^2)}{x^2}|\le u^2$。
\end{example}
\subsection{第24法：收敛级数必要条件}\label{mat:limits_35:k24}
\begin{method}{第24法：收敛级数必要条件}\textbf{适用条件} 已知级数收敛，可利用项或尾和趋零。\par\textbf{实施与依据} 先确立收敛，再用Cauchy尾和估计。\end{method}
\begin{example}{专题例题｜第24法：收敛级数必要条件}{难度：$\star\star$}
\textbf{题目} $a_n\ge0$且$\sum a_n<\infty$，求$\lim\sum_{k=n}^{2n}a_k$。\par\textbf{思路与解答} 该和是收敛正项级数的尾和的一部分，夹在0与$\sum_{k=n}^\infty a_k\to0$之间，极限0。\par\textbf{方法点评与边界} 一般正项收敛不必有$na_n\to0$，除非再加单调等条件。
\end{example}
\subsection{第25法：级数求和求极限}\label{mat:limits_35:k25}
\begin{method}{第25法：级数求和求极限}\textbf{适用条件} 有限和是收敛级数的部分和，或有幂级数生成函数。\par\textbf{实施与依据} 求完整级数和，再判尾项消失。\end{method}
\begin{example}{专题例题｜第25法：级数求和求极限}{难度：$\star\star$}
\textbf{题目} 求$\lim_{N\to\infty}\sum_{n=1}^N\frac{n}{2^n}$。\par\textbf{思路与解答} 幂级数$\sum nx^n=\frac{x}{(1-x)^2}$，在$x=\frac{1}{2}$得到2；尾项由几何衰减控制。\par\textbf{方法点评与边界} Fourier选点求和必须先核对该点的收敛值。
\end{example}
\subsection{第26法：常见不等式}\label{mat:limits_35:k26}
\begin{method}{第26法：常见不等式}\textbf{适用条件} 平均、乘积、根式或指数需估计。\par\textbf{实施与依据} 选择AM--GM、Cauchy、Bernoulli或对数界。\end{method}
\begin{example}{专题例题｜第26法：常见不等式}{难度：$\star\star$}
\textbf{题目} 证明$(1+\frac{x}{n})^n\to\e^x$，$x>0$。\par\textbf{思路与解答} $\frac{u}{1+u}\le\ln(1+u)\le u$给$\frac{x}{1+\frac{x}{n}}\le n\ln(1+\frac{x}{n})\le x$，夹逼后取指数。\par\textbf{方法点评与边界} 不等式方向依赖参数范围。
\end{example}
\subsection{第27法：等价无穷大量}\label{mat:limits_35:k27}
\begin{method}{第27法：等价无穷大量}\textbf{适用条件} 增长量的比有非零有限极限。\par\textbf{实施与依据} 提取最大阶，或用积分估计先得到等价式。\end{method}
\begin{example}{专题例题｜第27法：等价无穷大量}{难度：$\star\star$}
\textbf{题目} 求$\lim_{n\to\infty}\frac{\ln(n!)}{n\ln n}$。\par\textbf{思路与解答} 由$\int_1^n\ln t\dd t\le\ln(n!)\le n\ln n$，下界为$n\ln n-n+1$，除以$n\ln n$夹逼得1。\par\textbf{方法点评与边界} 无穷大量的差也不能不经误差分析直接替换。
\end{example}
\subsection{第28法：Cauchy收敛准则}\label{mat:limits_35:k28}
\begin{method}{第28法：Cauchy收敛准则}\textbf{适用条件} 难以先猜出极限，但尾项差易估计。\par\textbf{实施与依据} 对所有$m>n$建立统一小界。\end{method}
\begin{example}{专题例题｜第28法：Cauchy收敛准则}{难度：$\star\star$}
\textbf{题目} 证明$S_n=\sum_{k=1}^n\frac{1}{k^2}$收敛。\par\textbf{思路与解答} 当$m>n\ge1$，$0\le S_m-S_n\le\int_n^\infty t^{-2}\dd t=\frac{1}{n}$。因此是Cauchy列，实数完备性给收敛。\par\textbf{方法点评与边界} 只估计$S_{n+1}-S_n\to0$不够。
\end{example}
\subsection{第29法：比值极限与根值极限}\label{mat:limits_35:k29}
\begin{method}{第29法：比值极限与根值极限}\textbf{适用条件} $a_n>0$且$\frac{a_{n+1}}{a_n}\to L$。\par\textbf{实施与依据} 把连续相邻比相乘，使用几何均值或对数Cesaro平均。\end{method}
\begin{example}{专题例题｜第29法：比值极限与根值极限}{难度：$\star\star$}
\textbf{题目} 若$\frac{a_{n+1}}{a_n}\to3$，求$\lim\sqrt[n]{a_n}$。\par\textbf{思路与解答} 任取$\varepsilon\in(0,3)$，最终比值夹在$3\pm\varepsilon$，连乘后取$n$次根得到同界，故极限3。\par\textbf{方法点评与边界} 逆命题不成立；根值趋3不强迫相邻比收敛。
\end{example}
\subsection{第30法：上下极限}\label{mat:limits_35:k30}
\begin{method}{第30法：上下极限}\textbf{适用条件} 有界振荡列或需判收敛。\par\textbf{实施与依据} 计算尾部上确界、下确界的极限；两者相等才收敛。\end{method}
\begin{example}{专题例题｜第30法：上下极限}{难度：$\star\star$}
\textbf{题目} $a_n=(-1)^n+\frac{1}{n}$的上下极限。\par\textbf{思路与解答} 偶数项趋1，奇数项趋$-1$，且整体最终落在两端任意小邻域内，故$\limsup a_n=1,\liminf a_n=-1$。\par\textbf{方法点评与边界} 上下极限可以存在而普通极限不存在。
\end{example}
\subsection{第31法：数学常数识别}\label{mat:limits_35:k31}
\begin{method}{第31法：数学常数识别}\textbf{适用条件} 表达式与某常数的定义相符。\par\textbf{实施与依据} 必须先写常数的定义或已经证明的表征。\end{method}
\begin{example}{专题例题｜第31法：数学常数识别}{难度：$\star\star$}
\textbf{题目} 证明$H_n-\ln n$收敛并记其极限。\par\textbf{思路与解答} $b_{n+1}-b_n=\frac{1}{n+1}-\ln(1+\frac{1}{n})<0$；且$H_n\ge\int_1^{n+1}\frac{\dd t}{t}$给$b_n>0$。故收敛，其极限定义为Euler常数$\gamma$。\par\textbf{方法点评与边界} 识别名称不是证明存在性；Catalan常数可定义为$\sum\frac{(-1)^n}{(2n+1)^2}$。
\end{example}
\subsection{第32法：Wallis公式}\label{mat:limits_35:k32}
\begin{method}{第32法：Wallis公式}\textbf{适用条件} 三角幂积分、双阶乘、中心二项式。\par\textbf{实施与依据} 由$I_n$递推建立相邻奇偶项夹逼，再转成乘积。\end{method}
\begin{example}{专题例题｜第32法：Wallis公式}{难度：$\star\star$}
\textbf{题目} 求$\lim_{n\to\infty}\frac{\sqrt{\pi n}\binom{2n}{n}}{4^n}$。\par\textbf{思路与解答} Wallis乘积或由$I_{2n+1}\le I_{2n}\le I_{2n-1}$与递推得到$\binom{2n}{n}\sim\frac{4^n}{\sqrt{\pi n}}$，极限1。\par\textbf{方法点评与边界} 此处已使用由积分夹逼推导的Wallis渐近结论，不能把待证渐近式当作定义。
\end{example}
\subsection{第33法：Stirling公式}\label{mat:limits_35:k33}
\begin{method}{第33法：Stirling公式}\textbf{适用条件} 大阶乘、组合数。\par\textbf{实施与依据} 用$n!\sim\sqrt{2\pi n}(\frac{n}{\e})^n$，精确跟踪幂次和常数。\end{method}
\begin{example}{专题例题｜第33法：Stirling公式}{难度：$\star\star$}
\textbf{题目} 求$\lim_{n\to\infty}\frac{n!\e^n}{n^n\sqrt n}$。\par\textbf{思路与解答} 代入Stirling公式立即得到$\sqrt{2\pi}$。公式可由$\ln(n!)$的积分主项、Euler--Maclaurin余项及Wallis确定常数推导。\par\textbf{方法点评与边界} 只用$\ln(n!)\sim n\ln n$无法恢复常数$\sqrt{2\pi}$。
\end{example}
\subsection{第34法：线性代数方法}\label{mat:limits_35:k34}
\begin{method}{第34法：线性代数方法}\textbf{适用条件} 常系数向量递推或矩阵幂。\par\textbf{实施与依据} 找特征值与不变子空间，主特征值决定增长。\end{method}
\begin{example}{专题例题｜第34法：线性代数方法}{难度：$\star\star$}
\textbf{题目} $v_{n+1}=Av_n$，$A=\begin{pmatrix}2&1\\1&2\end{pmatrix},v_0=(1,0)^T$，求$3^{-n}v_n$极限。\par\textbf{思路与解答} 基向量$(1,1),(1,-1)$的特征值分别3、1，$v_n=\frac12[3^n(1,1)+(1,-1)]$，归一化极限$(\frac{1}{2},\frac{1}{2})$。\par\textbf{方法点评与边界} Jordan块含额外多项式因子，不能只看特征值而忽略重根结构。
\end{example}
\subsection{第35法：概率方法与Bernstein逼近}\label{mat:limits_35:k35}
\begin{method}{第35法：概率方法与Bernstein逼近}\textbf{适用条件} 二项式权重非负且和为1，可解释成期望。\par\textbf{实施与依据} 用均值方差控制远离中心的权重，再利用连续性。\end{method}
\begin{example}{专题例题｜第35法：概率方法与Bernstein逼近}{难度：$\star\star$}
\textbf{题目} 对$f\in C[0,1]$证明$B_nf(x)=\sum_{k=0}^nf(\frac{k}{n})\binom nkx^k(1-x)^{n-k}\to f(x)$一致。\par\textbf{思路与解答} 令$X\sim\operatorname{Bin}(n,x)$，$B_nf=\mathbb E f(\frac{X}{n})$。对$\delta>0$，误差$\le\omega_f(\delta)+2\|f\|_\infty\mathbb P(|\frac{X}{n}-x|>\delta)$。方差$\le\frac{1}{4n}$，Chebyshev给误差$\le\omega_f(\delta)+\frac{\|f\|_\infty}{2n\delta^2}$。先选$\delta$再取$n$，得一致收敛。\par\textbf{方法点评与边界} 此法同时证明连续函数可被多项式一致逼近；概率只是非负权重估计的组织语言。
\end{example}




\section{极限方法的证明、辨析与多解训练}\label{deep:limits_source_deep}

\begin{examtip}本节在35法方法卡之外补充证明、误差阶和方法比较。例题分别标为考研风格或竞赛风格，训练条件辨识与方法选择。\end{examtip}

\subsection{定义、否定命题与子列判别}\label{deep:limits_source_deep:logic}

\begin{method}{定义、否定命题与子列判别}\textbf{适用条件} 先确定自变量趋向方式和函数定义域。\par\textbf{步骤与依据} 把“任意$\varepsilon$”与“存在$\delta$”按次序写出；否定时交换量词；不收敛可构造两个极限不同的子列。\par\textbf{易错点} “某些项接近”只说明存在子列，不能说明整列收敛。\end{method}

\begin{theorem}{Cauchy准则与完备性}实数列收敛当且仅当对每个$\varepsilon>0$存在$N$，使$m,n>N$时$|a_m-a_n|<\varepsilon$。\end{theorem}\begin{proof}必要性由$|a_m-a_n|\le |a_m-a|+|a_n-a|$。充分性：Cauchy列有界，由Bolzano--Weierstrass定理有子列$a_{n_k}\to a$。固定$\varepsilon$，取尾部Cauchy误差小于$\frac{\varepsilon}{2}$，再取$n_k$足够大使$|a_{n_k}-a|<\frac{\varepsilon}{2}$，三角不等式给全部尾项趋于$a$。\end{proof}

\begin{example}{基础例题｜量词与误差估计}{难度：$\star\star$}\label{deep:ex:limits_source_deep:eps}\textbf{题目} 用定义证明$a_n=\frac{2n+1}{3n-2}\to\frac{2}{3}$。\par\textbf{思路} 精确计算与候选极限的差，再从误差反解$n$。\par\textbf{详细解答} $|a_n-\frac{2}{3}|=\frac{7}{3(3n-2)}$。当$n\ge2$，$3n-2\ge2n$，故误差$\le\frac{7}{6n}$。给定$\varepsilon>0$，取整数$N>\max\{2,\frac{7}{6\varepsilon}\}$，则$n>N$时误差小于$\varepsilon$。\par\textbf{方法点评} 先放缩再选$N$可以简化常数。\par\textbf{变式与易错点} 把$n$换为$x\to+\infty$可写成函数极限定义。\end{example}

\begin{example}{竞赛风格｜Heine定理的双向使用}{难度：$\star\star\star$}\label{deep:ex:limits_source_deep:heine}\textbf{题目} 证明$\sin(\frac{1}{x})$在$x\to0$时无极限，而$x\sin(\frac{1}{x})\to0$。\par\textbf{思路} 前者选相位相差$\pi$的两列；后者用有界性消除振荡。\par\textbf{详细解答} 取$x_n=(2\pi n+\frac{\pi}{2})^{-1}$和$y_n=(2\pi n+\frac{3\pi}{2})^{-1}$，均趋0，而函数值分别为1、$-1$，由Heine定理极限不存在。又$-|x|\le x\sin(\frac{1}{x})\le|x|$，夹逼得极限0。\par\textbf{方法点评} 子列用于否定极限；有界乘无穷小用于建立极限。\par\textbf{变式与易错点} 振荡函数与趋零函数的积不一定振荡。\end{example}

\subsection{相消、等价无穷小与Taylor余项}\label{deep:limits_source_deep:cancel}

\begin{method}{相消、等价无穷小与Taylor余项}\textbf{适用条件} 分子有相减且首项可能抵消。\par\textbf{步骤与依据} 先估计分母阶数；展开到抵消后第一项，并检验余项除以分母后趋0。\par\textbf{易错点} 在差式里逐项等价替换可能把真正的主项消掉。\end{method}

\begin{example}{考研风格｜同一极限三解}{难度：$\star\star\star$}\label{deep:ex:limits_source_deep:sin}\textbf{题目} 求$L=\lim_{x\to0}\frac{x-\sin x}{x^3}$。\par\textbf{思路} 分别用Taylor、L’Hôpital和积分余项。\par\textbf{详细解答} \textbf{方法一：Taylor。}$\sin x=x-\frac{x^3}{6}+o(x^3)$，所以$L=\frac{1}{6}$。\par\textbf{方法二：L’Hôpital。}三次求导后依次得到$\frac{1-\cos x}{3x^2}$、$\frac{\sin x}{6x}$、$\frac{\cos x}{6}$，极限$\frac{1}{6}$；每次在穿孔邻域内导数分母非零。\par\textbf{方法三：积分。}$x-\sin x=x\int_0^1(1-\cos xu)\dd u$，故原商$=\int_0^1\frac{1-\cos xu}{x^2}\dd u$。被积函数趋$\frac{u^2}{2}$，且绝对值$\le \frac{u^2}{2}$，积分趋$\frac{1}{6}$。\par\textbf{方法点评} Taylor最短；积分法容易推广为带权平均。\par\textbf{变式与易错点} 不能将差式里的$\sin x$直接替成$x$。\end{example}

\begin{example}{考研风格｜参数相消与分类}{难度：$\star\star\star$}\label{deep:ex:limits_source_deep:log}\textbf{题目} 求$\lim_{x\to0}\frac{\ln(1+ax)-a\ln(1+x)}{x^2}$，$a\in\R$。\par\textbf{思路} 线性项必相消，保留二阶项。\par\textbf{详细解答} \textbf{展开法：}分子$=(-\frac{a^2}{2}+\frac{a}{2})x^2+o(x^2)$，极限$\frac{a(1-a)}{2}$。\par\textbf{积分法：}分子导数为$\frac{a(1-a)x}{(1+ax)(1+x)}$。从0积分，令$t=xu$，原商$=a(1-a)\int_0^1\frac{u}{(1+axu)(1+xu)}\dd u$，得同值。\par\textbf{方法点评} 积分法自动解释$a=0,1$时恒为0。\par\textbf{变式与易错点} 固定$a$后取足够小的$x$保证对数有定义。\end{example}

\begin{example}{考研风格｜指数型的二阶精度}{难度：$\star\star\star$}\label{deep:ex:limits_source_deep:exponential}\textbf{题目} 求$\lim_{n\to\infty}n^2[\frac{(1+\frac{1}{n})^n}{\e}-1+\frac{1}{2n}]$。\par\textbf{思路} 先对数展开到$n^{-2}$，再指数回代。\par\textbf{详细解答} $n\ln(1+\frac{1}{n})=1-\frac1{2n}+\frac1{3n^2}+O(n^{-3})$。令$u_n=-\frac{1}{2n}+\frac{1}{3n^2}+O(n^{-3})$，则$\e^{u_n}=1+u_n+\frac{u_n^2}{2}+O(n^{-3})=1-\frac1{2n}+\frac{11}{24n^2}+O(n^{-3})$，答案$\frac{11}{24}$。另一法对$f(t)=\exp(\frac{\ln(1+t)}{t}-1)$在0补值，求其二阶Taylor系数，仍得$\frac{11}{24}$。\par\textbf{方法点评} 外层指数的平方项贡献$\frac{1}{8}$。\par\textbf{变式与易错点} 仅用$\e^{u_n}-1\sim u_n$会丢失二阶余量。\end{example}

\begin{example}{考研风格｜根式的三种组织}{难度：$\star\star\star$}\label{deep:ex:limits_source_deep:rationalization}\textbf{题目} 求$\lim_{x\to0}\frac{\sqrt{1+x}-\sqrt{1-x}-x}{x^3}$。\par\textbf{思路} 奇函数结构消去偶次项。\par\textbf{详细解答} \textbf{Taylor法：}两平方根相减为$x+\frac{x^3}{8}+o(x^3)$，答案$\frac{1}{8}$。\par\textbf{有理化法：}设$A=\sqrt{1+x}+\sqrt{1-x}$，原商$=\frac{2-A}{Ax^2}$。又$2-A=\frac{2x^2}{(2+A)(1+\sqrt{1-x^2})}$，代入趋$\frac{1}{8}$。\par\textbf{导数法：}分子$h$满足$h(0)=h'(0)=h''(0)=0,h'''(0)=\frac{3}{4}$，Taylor积分余项给$\frac{h}{x^3}\to \frac{h'''(0)}{6}=\frac{1}{8}$。\par\textbf{方法点评} 有理化显示隐藏的$x^2$因子。\par\textbf{变式与易错点} 第一次有理化后仍有相消，须继续处理。\end{example}

\subsection{Stolz定理、加权平均与求和}\label{deep:limits_source_deep:stolz}

\begin{method}{Stolz定理、加权平均与求和}\textbf{适用条件} $b_n$严格递增且趋$+\infty$，差分比有极限。\par\textbf{步骤与依据} 把$a_n-a_N$拆成增量求和，利用尾部差分比上下界。\par\textbf{易错点} 分母不单调或不无界时，不可套该版本。\end{method}

\begin{theorem}{Stolz--Cesàro定理及证明}若$b_n\uparrow+\infty$且$\frac{a_{n+1}-a_n}{b_{n+1}-b_n}\to L\in\R$，则$\frac{a_n}{b_n}\to L$。\end{theorem}\begin{proof}对$n\ge N$有$(L-\varepsilon)(b_{n+1}-b_n)\le a_{n+1}-a_n\le(L+\varepsilon)(b_{n+1}-b_n)$。相加后除以$b_n$。固定首段趋0，且$\frac{b_n-b_N}{b_n}\to1$，上下极限在$L\pm\varepsilon$之间。令$\varepsilon\downarrow0$。\end{proof}

\begin{example}{竞赛风格｜幂和的三解}{难度：$\star\star\star$}\label{deep:ex:limits_source_deep:powersum}\textbf{题目} $p>0$时求$\lim_{n\to\infty}n^{-p-1}\sum_{k=1}^nk^p$。\par\textbf{思路} 差分、积分和、单调积分夹逼。\par\textbf{详细解答} \textbf{Stolz法：}差分比分子$(n+1)^p$，分母$(n+1)^{p+1}-n^{p+1}\sim(p+1)n^p$，答案$\frac{1}{p+1}$。\par\textbf{Riemann和：}原式$=\frac1n\sum(\frac{k}{n})^p\to\int_0^1t^p\dd t$。\par\textbf{夹逼：}$\int_0^nx^p\dd x\le\sum_{k=1}^nk^p\le\int_1^{n+1}x^p\dd x$，同除$n^{p+1}$。\par\textbf{方法点评} 积分和直观；Stolz适合一般增量；夹逼还能估误差。\par\textbf{变式与易错点} $-1<p<0$时端点奇异，需要额外估计。\end{example}

\begin{example}{竞赛风格｜加权平均原理}{难度：$\star\star\star$}\label{deep:ex:limits_source_deep:weighted}\textbf{题目} $a_n\to A,w_n>0,W_n=\sum_{k=1}^nw_k\to\infty$，证明$\sum \frac{w_ka_k}{W_n}\to A$。\par\textbf{思路} 有限首段与统一小的尾段分开。\par\textbf{详细解答} 取$N$使$k>N$时$|a_k-A|<\varepsilon$。误差$\le W_n^{-1}\sum_{k=1}^Nw_k|a_k-A|+\varepsilon=o(1)+\varepsilon$。另一法令$S_n=\sum w_ka_k$，增量比$\frac{S_{n+1}-S_n}{W_{n+1}-W_n}=a_{n+1}\to A$，用Stolz。\par\textbf{方法点评} 包含Cesàro均值和调和加权极限。\par\textbf{变式与易错点} 若总权重有限，首段影响一般不会消失。\end{example}

\begin{example}{竞赛风格｜阶乘与对数平均}{难度：$\star\star\star$}\label{deep:ex:limits_source_deep:factorial}\textbf{题目} 不调用Stirling公式，求$\lim_{n\to\infty}\frac{(n!)^{\frac{1}{n}}}{n}$。\par\textbf{思路} 取对数并处理$\ln t$在0的奇点。\par\textbf{详细解答} 单调积分比较给$\int_1^n\ln x\dd x\le\ln(n!)\le\int_1^n\ln x\dd x+\ln n$。所以$\frac{\ln(n!)}{n}-\ln n=-1+O(\frac{\ln n}{n})$，原极限$\e^{-1}$。第二法对$\frac{\ln(n!)-n\ln n}{n}$用Stolz，增量$=-n\ln(1+\frac{1}{n})\to-1$。\par\textbf{方法点评} 两法只用初等重要极限。\par\textbf{变式与易错点} $\ln t$不在$[0,1]$连续，积分和须处理端点。\end{example}

\subsection{单调有界、压缩与非线性迭代}\label{deep:limits_source_deep:recurrence}

\begin{method}{单调有界、压缩与非线性迭代}\textbf{适用条件} 递推由连续映射给出，能找到不变区间。\par\textbf{步骤与依据} 先证每项留在区间，再比较相邻项；导数绝对值小于1时给收敛速度。\par\textbf{易错点} 解不动点方程只能列候选值。\end{method}

\begin{example}{考研风格｜迭代极限及误差}{难度：$\star\star\star$}\label{deep:ex:limits_source_deep:sqrt}\textbf{题目} $x_1=1,x_{n+1}=\sqrt{2+x_n}$，证明收敛并估计速度。\par\textbf{思路} $[1,2]$不变，比较$x_{n+1}^2$与$x_n^2$。\par\textbf{详细解答} 若$1\le x_n<2$，则$1<x_{n+1}<2$，且$2+x_n-x_n^2=(2-x_n)(1+x_n)>0$，故递增。极限满足$L=\sqrt{2+L}$，区间内为2。又$2-x_{n+1}=\frac{2-x_n}{2+\sqrt{2+x_n}}\le\frac{2-x_n}{3}$，所以$2-x_n\le3^{1-n}$。第二法在区间上$|T'|\le\frac{1}{2\sqrt3}<1$，用压缩。\par\textbf{方法点评} 单调法证明存在，精确误差关系给速度。\par\textbf{变式与易错点} 平方产生的候选$-1$不满足根号非负。\end{example}

\begin{example}{竞赛风格｜Newton迭代的二次收敛}{难度：$\star\star\star$}\label{deep:ex:limits_source_deep:newton}\textbf{题目} $a>0,x_1>0,x_{n+1}=\frac{x_n+\frac{a}{x_n}}{2}$，证明$x_n\to\sqrt a$。\par\textbf{思路} 均值不等式给下界，因式分解给误差。\par\textbf{详细解答} $x_{n+1}\ge\sqrt a$。从$n\ge2$，$x_{n+1}-x_n=\frac{a-x_n^2}{2x_n}\le0$，有极限且为$\sqrt a$。精确式$x_{n+1}-\sqrt a=\frac{(x_n-\sqrt a)^2}{2x_n}$给二次速度。另设$r_n=\frac{x_n-\sqrt a}{x_n+\sqrt a}$，化简得$r_{n+1}=r_n^2$，所以$r_n=r_1^{2^{n-1}}$，也证明收敛。\par\textbf{方法点评} 变量变换揭示迭代的闭式结构。\par\textbf{变式与易错点} $x_1=0$无定义；负初值须另作讨论。\end{example}

\subsection{奇异Riemann和与端点集中}\label{deep:limits_source_deep:integral}

\begin{method}{奇异Riemann和与端点集中}\textbf{适用条件} 采样函数端点无界但可积，或权重集中在端点。\par\textbf{步骤与依据} 切分近端与远端，分别用单调积分和一致连续性控制。\par\textbf{易错点} 不能交换未经估计的极限与积分。\end{method}

\begin{example}{竞赛风格｜带可积奇点的和}{难度：$\star\star\star$}\label{deep:ex:limits_source_deep:singular}\textbf{题目} $0<\alpha<1$时求$\lim_{n\to\infty}n^{\alpha-1}\sum_{k=1}^nk^{-\alpha}$。\par\textbf{思路} 用单调积分比较将奇异积分和严格化。\par\textbf{详细解答} $\int_1^{n+1}x^{-\alpha}\dd x\le\sum_{k=1}^nk^{-\alpha}\le1+\int_1^nx^{-\alpha}\dd x$。同乘$n^{\alpha-1}$，两界趋$\frac{1}{1-\alpha}$。第二法Stolz对$\sum \frac{k^{-\alpha}}{n^{1-\alpha}}$作差，同样得$\frac{1}{1-\alpha}$。\par\textbf{方法点评} 比较也解释$\alpha=1$时对数增长，$\alpha>1$时和有界。\par\textbf{变式与易错点} 不能把$\alpha\ge1$代入有限积分值。\end{example}

\begin{example}{竞赛风格｜端点集中的积分}{难度：$\star\star\star$}\label{deep:ex:limits_source_deep:laplace}\textbf{题目} $f\in C[0,1]$，求$\lim_{n\to\infty}(n+1)\int_0^1x^nf(x)\dd x$。\par\textbf{思路} 归一化权重质量集中在1。\par\textbf{详细解答} 减去$f(1)$。取$\delta$使$x>1-\delta$时$|f(x)-f(1)|<\varepsilon$，该段误差$\le\varepsilon$。其余段误差$\le2\|f\|_\infty(1-\delta)^{n+1}\to0$，故答案$f(1)$。若$f\in C^1$，分部得$(n+1)\int x^nf=f(1)-\int_0^1x^{n+1}f'(x)\dd x$，也得结论。\par\textbf{方法点评} 第一法只需连续，分部法假设更强。\par\textbf{变式与易错点} 无$n+1$因子时积分趋0。\end{example}

\subsection{上下极限、增长阶与经典常数}\label{deep:limits_source_deep:asymptotic}

\begin{method}{上下极限、增长阶与经典常数}\textbf{适用条件} 数列振荡，或只有比值、根式信息。\par\textbf{步骤与依据} 尾部上确界定义$\limsup$；比值界相乘后取$n$次根。\par\textbf{易错点} 比值极限不存在时不能强行写根式极限与它相等。\end{method}

\begin{theorem}{比值控制根式}$a_n>0$时$\liminf \frac{a_{n+1}}{a_n}\le\liminf\sqrt[n]{a_n}\le\limsup\sqrt[n]{a_n}\le\limsup \frac{a_{n+1}}{a_n}$，允许扩展实数。\end{theorem}\begin{proof}若$q>\limsup \frac{a_{n+1}}{a_n}$且$q>0$，尾部$a_{n+1}\le q a_n$，故$a_n\le Cq^n$，得到$\limsup\sqrt[n]{a_n}\le q$，再令$q$下降。下界由$a_{n+1}\ge p a_n$同证；0和无穷情形以任意正$q,p$处理。\end{proof}

\begin{example}{竞赛风格｜Euler常数存在性}{难度：$\star\star\star$}\label{deep:ex:limits_source_deep:gamma}\textbf{题目} 证明$H_n-\ln n$收敛，$H_n=\sum_{k=1}^n\frac{1}{k}$。\par\textbf{思路} 对数上下界证明单调有界。\par\textbf{详细解答} $a_{n+1}-a_n=\frac{1}{n+1}-\ln(1+\frac{1}{n})<0$。又$H_n\ge\int_1^{n+1}\frac{\dd x}{x}=\ln(n+1)$，故$a_n>0$。极限记$\gamma$。尾差$0<a_n-\gamma<\frac{1}{n}$由$\ln(1+\frac{1}{k})-\frac{1}{k+1}<\frac{1}{k}-\frac{1}{k+1}$累加后令终点趋无穷得到。\par\textbf{方法点评} “和减积分”把离散求和与连续估计联系起来。\par\textbf{变式与易错点} 此证明不涉及$\gamma$是否无理。\end{example}

\begin{example}{竞赛风格｜Wallis乘积证明链}{难度：$\star\star\star$}\label{deep:ex:limits_source_deep:wallis}\textbf{题目} 令$I_n=\int_0^{\frac{\pi}{2}}\sin^nx\dd x$，推导Wallis乘积。\par\textbf{思路} 分部递推后用相邻项夹逼。\par\textbf{详细解答} $I_n=\frac{(n-1)I_{n-2}}{n}$，$I_0=\frac{\pi}{2},I_1=1$，故$I_{2n}=\frac{(\frac{\pi}{2})(2n-1)!!}{(2n)!!}$，$I_{2n+1}=\frac{(2n)!!}{(2n+1)!!}$。又$I_{2n+1}\le I_{2n}\le I_{2n-1}=\frac{2n+1}{2n}I_{2n+1}$，得$\frac{I_{2n}}{I_{2n+1}}\to1$。代入得$\prod_{k=1}^n\frac{(2k)^2}{(2k-1)(2k+1)}\to\frac{\pi}{2}$。\par\textbf{方法点评} 极限由夹逼保证，递推负责化乘积。\par\textbf{变式与易错点} 无限乘积指有限乘积极限。\end{example}

\begin{example}{竞赛风格｜矩阵迭代与主特征值}{难度：$\star\star\star$}\label{deep:ex:limits_source_deep:matrix}\textbf{题目} $u_{n+2}=u_{n+1}+u_n,u_0=0,u_1=1$，求$\frac{u_{n+1}}{u_n}$极限。\par\textbf{思路} 比较特征根模长，或改用比值递推。\par\textbf{详细解答} 根为$\phi=\frac{1+\sqrt5}{2},\psi=\frac{1-\sqrt5}{2}$，$u_n=\frac{\phi^n-\psi^n}{\sqrt5}$，因$|\frac{\psi}{\phi}|<1$，答案$\phi$。第二法$r_{n+1}=1+\frac{1}{r_n}$；从$n\ge2$留在$[\frac{3}{2},2]$，映射导数绝对值$\le\frac{4}{9}$，压缩到唯一不动点$\phi$。\par\textbf{方法点评} 矩阵法可推广高阶递推；压缩法无需显式解。\par\textbf{变式与易错点} $u_0=0$不能作分母。\end{example}

\begin{example}{竞赛风格｜Bernstein逼近与概率法}{难度：$\star\star\star$}\label{deep:ex:limits_source_deep:probability}\textbf{题目} 对$f\in C[0,1]$证明$B_nf(x)=\sum_{k=0}^nf(\frac{k}{n})\binom nkx^k(1-x)^{n-k}$一致趋$f$。\par\textbf{思路} 用二项分布的方差控制远点误差。\par\textbf{详细解答} 设$X\sim\mathrm{Bin}(n,x)$，$B_nf(x)=\mathbb E f(\frac{X}{n})$，且$\mathbb E(\frac{X}{n})=x$、$\operatorname{Var}(\frac{X}{n})=\frac{x(1-x)}{n}\le\frac{1}{4n}$。近点$|\frac{X}{n}-x|\le\delta$误差$\le\omega_f(\delta)$；远点误差$\le2\|f\|_\infty\Pr(|\frac{X}{n}-x|>\delta)\le\frac{\|f\|_\infty}{2n\delta^2}$。先取$\delta$小，再取$n$大，得到一致收敛。也可不用概率术语，直接用二项式恒等式计算同一二阶矩。\par\textbf{方法点评} 连接概率估计与多项式逼近。\par\textbf{变式与易错点} 一致连续性保证模连续性趋0。\end{example}

\subsection{实数构造与无理逼近}\label{deep:limits_source_deep:real}

\begin{method}{实数构造与无理逼近}\textbf{适用条件} 运用实数完备性与整数格点构造极限和逼近。\par\textbf{步骤与依据} 完备性保证上确界；整数格点给有理逼近；抽屉原理提高逼近精度。\par\textbf{易错点} 最优逼近指数不能从一次小数截断推断。\end{method}

\begin{example}{竞赛拓展｜Dirichlet有理逼近}{难度：$\star\star\star$}\label{deep:ex:limits_source_deep:dirichlet}\textbf{题目} 证明每个无理数$\alpha$有无穷多个$\frac{p}{q}$满足$|\alpha-\frac{p}{q}|<\frac{1}{q^2}$。\par\textbf{思路} 将$0,\alpha,\ldots,N\alpha$的小数部分放入$N$个等长区间。\par\textbf{详细解答} 两点落同区间，故有$1\le q\le N$和$p\in\mathbb Z$使$|q\alpha-p|<\frac{1}{N}$，于是$|\alpha-\frac{p}{q}|<\frac{1}{qN}\le\frac{1}{q^2}$。分母必无界：否则有限个$q$对应的非零距离$|q\alpha-p|$有正下界，与$\frac{1}{N}\to0$矛盾。\par\textbf{方法点评} 给出无理数有理逼近的基本构造机制。\par\textbf{变式与易错点} 不能据此确定特定常数的最优无理性指数。\end{example}

\subsection{Stirling公式：存在常数与确定常数}\label{deep:limits_source_deep:stirlingproof}

\begin{method}{Stirling公式：存在常数与确定常数}\textbf{适用条件} 需要阶乘的相对误差，而不只是$n$次根的极限。\par\textbf{步骤与依据} 先构造对数修正项证明收敛，再用Wallis确定常数；两步缺一不可。\par\textbf{易错点} 只用$\ln n!\sim n\ln n$不能推出$n!\sim\sqrt{2\pi n}(\frac{n}{\e})^n$。\end{method}

\begin{theorem}{Stirling等价式的证明}$n!\sim\sqrt{2\pi n}(\frac{n}{\e})^n$。\end{theorem}\begin{proof}令$s_n=\ln(n!)-(n+\frac12)\ln n+n$，则$s_{n+1}-s_n=1-(n+\frac12)\ln(1+\frac{1}{n})=-\frac{1}{12n^2}+O(n^{-3})$。增量绝对可和，故$s_n\to c$，并有$s_n=c+O(\frac{1}{n})$。所以$n!=\e^c n^{n+\frac{1}{2}}\e^{-n}(1+O(\frac{1}{n}))$。\par 由Wallis积分的精确乘积$I_{2n}I_{2n+1}=\frac{\pi}{2(2n+1)}$以及比值$\frac{I_{2n}}{I_{2n+1}}\to1$，得$I_{2n}\sim\frac{\sqrt\pi}{2\sqrt n}$。又$I_{2n}=(\frac{\pi}{2})\binom{2n}n4^{-n}$，因此$\binom{2n}n\sim\frac{4^n}{\sqrt{\pi n}}$。另一方面将含$c$的阶乘式代入二项系数，得$\binom{2n}n\sim\frac{\e^{-c}\sqrt2\,4^n}{\sqrt n}$。比较常数可得$\e^c=\sqrt{2\pi}$，证毕。\end{proof}

\begin{example}{竞赛风格｜阶乘比的三条入口}{难度：$\star\star\star$}\label{deep:ex:limits_source_deep:stirlingcase}\textbf{题目} 求$\lim_{n\to\infty}\frac{\sqrt n\,\binom{2n}n}{4^n}$。\par\textbf{思路} 可用Stirling、Wallis，或中心二项系数的积分表示。\par\textbf{详细解答} Stirling直接给$\frac{1}{\sqrt\pi}$。第二法由上证$I_{2n}\sim\frac{\sqrt\pi}{2\sqrt n}$且$\frac{\binom{2n}n}{4^n}=\frac{2I_{2n}}{\pi}$，同得$\frac{1}{\sqrt\pi}$。第三条解释是$\int_{-\frac{\pi}{2}}^{\frac{\pi}{2}}\cos^{2n}t\dd t$在0附近集中，缩放$t=\frac{u}{\sqrt n}$后趋Gaussian积分$\int_\R\e^{-u^2}\dd u=\sqrt\pi$；要把这条解释变为证明，还需用$\cos t\le\e^{-ct^2}$控制尾部。\par\textbf{方法点评} 前两条是本册已完成的证明；第三条给出可进一步独立练习的积分渐近路线。\par\textbf{变式与易错点} 中央二项系数的$4^n$尺度和$n^{-\frac{1}{2}}$因子都不可漏。\end{example}

\begin{example}{竞赛拓展｜阶乘尾项与e的无理性}{难度：$\star\star\star$}\label{deep:ex:limits_source_deep:eirrat}\textbf{题目} 证明$\e$不是有理数。\par\textbf{思路} 若$\e=\frac{p}{q}$，乘$q!$后级数首段为整数，而尾段严格在0与1之间。\par\textbf{详细解答} 可取$q\ge2$。若$\e=\frac{p}{q}$，则$N=q!\e-\sum_{k=0}^q\frac{q!}{k!}$是整数。另一方面$N=\sum_{j\ge1}\frac{q!}{(q+j)!}>0$，且每项$\le(q+1)^{-j}$，并至少从$j=2$起严格，故$N<\sum_{j\ge1}(q+1)^{-j}=\frac{1}{q}<1$，矛盾。\par\textbf{方法点评} 关键是“整数却在$(0,1)$”，尾项的几何控制来自阶乘增长。\par\textbf{变式与易错点} 不能仅说$\e$的小数无限不循环而不证明。\end{example}

\section{极限计算的误差与增长构型}\label{ic:delta_limits}

\begin{examtip}本节讨论一阶Riemann和误差、非收敛递推的增长尺度及参数几何平均。\end{examtip}

\subsection{Riemann和的一阶误差}\label{ic:delta_limits:riemannerror}\begin{method}{Riemann和的一阶误差}\textbf{识别信号} $f\in C^1[0,1]$。\par \textbf{操作步骤} 把每个小区间内端点值与积分的差写成$f'$的积分；用$f'$一致连续性控制余项，再求和。\par \textbf{疑难点与条件} 左、右端点和的首项误差符号相反；普通Riemann和极限不能单独给误差系数。\end{method}

\begin{theorem}{右端点误差公式}设$f\in C^1[0,1]$，则$\lim_{n\to\infty}n[\frac1n\sum_{k=1}^nf(\frac{k}{n})-\int_0^1f]=\frac{f(1)-f(0)}{2}$。\end{theorem}\begin{proof}对$I_k=[\frac{k-1}{n},\frac{k}{n}]$，$f(\frac{k}{n})-f(t)=f'(\frac{k}{n})(\frac{k}{n}-t)+r_k(t)$，其中$|r_k(t)|\le\omega_{f'}(\frac{1}{n})(\frac{k}{n}-t)$。积分后乘$n$，主项为$\frac{f'(\frac{k}{n})}{2n}$，各余项累加绝对值$\le\frac{\omega_{f'}(\frac{1}{n})}{2}\to0$。主项的Riemann和趋$\frac12\int_0^1f'=\frac{f(1)-f(0)}{2}$。\end{proof}

\begin{example}{竞赛风格｜指数函数端点误差}{难度：$\star\star\star$}\label{ic:ex:delta_limits:riemannerrorcase}\textbf{题目} 求$\lim_{n\to\infty}n[\frac1n\sum_{k=1}^n\e^{\frac{k}{n}}-(\e-1)]$。\par \textbf{思路} 用误差公式或有限等比和。\par \textbf{详细解答} 公式直接得$\frac{\e-1}{2}$。独立地，$\frac1n\sum\e^{\frac{k}{n}}=\frac{(\e-1)\e^{\frac{1}{n}}}{n(\e^{\frac{1}{n}}-1)}$，Taylor给右边$(\e-1)[1+\frac{1}{2n}+O(n^{-2})]$，同值。\par \textbf{验算与条件复核} 两个独立计算核对误差系数及正号。\par \textbf{疑难点与方法比较} 若用左端点$k=0,\ldots,n-1$，答案变为$-\frac{\e-1}{2}$。\end{example}

\begin{example}{竞赛风格｜发散递推的精确增长阶}{难度：$\star\star\star$}\label{ic:ex:delta_limits:growthrecurrence}\textbf{题目} $a_1>0,a_{n+1}=a_n+\frac{1}{a_n}$，证明$\frac{a_n}{\sqrt{2n}}\to1$。\par \textbf{思路} 先证无界，再对平方量用Stolz。\par \textbf{详细解答} 数列正且严格递增。若有有限极限$L>0$，递推给$L=L+\frac{1}{L}$，矛盾，所以$a_n\to\infty$。又$a_{n+1}^2-a_n^2=2+\frac{1}{a_n^2}\to2$。Stolz给$\frac{a_n^2}{2n}\to1$，正平方根连续，故结论。\par \textbf{验算与条件复核} 平方增量是精确恒等式；取根时正性已证明。\par \textbf{疑难点与方法比较} 本题数列本身发散，但归一化后收敛，不能套不动点存在论。\end{example}

\begin{example}{竞赛风格｜两指数相消的二阶项}{难度：$\star\star\star$}\label{ic:ex:delta_limits:powercollision}\textbf{题目} 求$\lim_{x\to\e}\frac{x^\e-\e^x}{(x-\e)^2}$。\par \textbf{思路} 提$\e^x$再对剩余比值取对数。\par \textbf{详细解答} 令$h(x)=\e\ln x-x$，有$h(\e)=h'(\e)=0,h''(\e)=-\frac{1}{\e}$，所以$h(x)=-\frac{(x-\e)^2}{2\e}+o((x-\e)^2)$。原分子$\e^x(\e^{h(x)}-1)\sim-\frac{\e^{\e-1}(x-\e)^2}{2}$，答案$-\frac{\e^{\e-1}}{2}$。\par \textbf{验算与条件复核} 直接对$x^\e-\e^x$求二阶导在$\e$取值也为$-\e^{\e-1}$，除以$2!$相同。\par \textbf{疑难点与方法比较} 两项分别等价后直接相减会丢掉真正二阶项。\end{example}

\begin{example}{竞赛风格｜广义平均的根式极限}{难度：$\star\star\star$}\label{ic:ex:delta_limits:geometricmean}\textbf{题目} $a_1,\ldots,a_m>0$，求$\lim_{n\to\infty}[\frac{a_1^{\frac{1}{n}}+\cdots+a_m^{\frac{1}{n}}}{m}]^n$。\par \textbf{思路} 先对数，再对小参数$\frac{1}{n}$展开。\par \textbf{详细解答} 记$\ell_j=\ln a_j$，则平均底数$=1+\frac1{mn}\sum\ell_j+O(n^{-2})$。对数乘$n$趋$m^{-1}\sum\ell_j$，所以原极限$\exp(m^{-1}\sum\ell_j)=(\prod a_j)^{\frac{1}{m}}$。\par \textbf{验算与条件复核} 有限$m$使所有Taylor余项可统一控制；指数连续性完成回代。\par \textbf{疑难点与方法比较} 底数中的$m$不能省略，否则$n$次幂将改变尺度。\end{example}

\section{图解：极限的误差带}
\begin{figure}[htbp]\centering
\begin{tikzpicture}\begin{axis}[width=.8\textwidth,height=6cm,axis lines=middle,xmin=0,xmax=1.8,ymin=0,ymax=2.9,xlabel=$x$,ylabel=$f(x)$,samples=100]
\addplot[Teal,thick,domain=0:1.65]{x^2};
\addplot[dashed,Gold,domain=0:1.7]{1.3};\addplot[dashed,Gold,domain=0:1.7]{.7};
\draw[dashed,Ink] (axis cs:.9,0)--(axis cs:.9,2.5);
\draw[dashed,Ink] (axis cs:1.1,0)--(axis cs:1.1,2.5);
\fill[Ink] (axis cs:1,1) circle[radius=1.8pt];
\node[anchor=west] at (axis cs:1.25,1.35) {$L+\varepsilon$};
\node[anchor=west] at (axis cs:1.25,.65) {$L-\varepsilon$};
\node at (axis cs:1,2.65) {$|x-1|<\delta$};
\end{axis}\end{tikzpicture}
\caption{以$f(x)=x^2,x\to1$为例，竖直带控制自变量，水平带控制函数误差。图取$\delta=0.1,\varepsilon=0.3$；当$|x-1|<0.1$时$|x^2-1|=|x-1||x+1|<0.21<0.3$。定义要求对任意正$\varepsilon$都能选$\delta$，图中只展示其中一次选择。}\label{fig:review:limit}
\end{figure}

\section{保号、有限首段与零点定位}\label{cd:limits-sign}
\begin{examtip}
极限给出的是最终成立的固定精度估计。利用这种估计，可以把数列的最值问题压到有限首段，也可以在开区间或无界区间上选出适合介值定理的两个有限端点。
\end{examtip}

\subsection{极限附近的固定间隔与有限首段}
\begin{method}{先固定间隔，再截取尾部}
\textbf{适用条件} 数列$a_n\to L$，且已找到$a_p>L$或$a_q<L$。\par
\textbf{步骤} 用已知项与$L$之间的距离选定正数$\varepsilon$；将所有充分靠后的项收进$(L-\varepsilon,L+\varepsilon)$；最后只比较有限首段。\par
\textbf{疑难点} $\varepsilon$必须先固定。不能把“对任意固定$\varepsilon$最终成立”直接替换成$a_n-L<\frac{1}{n}$。
\end{method}
\begin{theorem}{收敛数列的最值取得判据}\label{cd:finite-extrema}
设$a_n\to L$。若存在$p$使$a_p>L$，则数列取得最大值；若存在$q$使$a_q<L$，则取得最小值。特别地，若数列在$L$两侧各有至少一项，则最大值、最小值均能取得。
\end{theorem}
\begin{proof}
令$\varepsilon=\frac{a_p-L}{2}>0$。存在$N\ge p$，使$n>N$时
\[
 a_n<L+\varepsilon=\dfrac{a_p+L}{2}<a_p.
\]
因此有限集合$\{a_1,\ldots,a_N\}$的最大值就是全数列最大值。对最小值取$\varepsilon=\frac{L-a_q}{2}$，同理可证。
\end{proof}
\begin{example}{考研风格｜有最值与仅有上下确界}{难度：$\star\star$}\label{cd:ex:finite-extrema}
\textbf{题目} 判断$a_n=\frac{\sin n}{n}$是否取得最大值和最小值；再与$b_n=\frac{1}{n}$比较。\par
\textbf{思路} 无须先找到最大项和最小项，只需找到极限两侧的项。\par
\textbf{详细解答} 由$|a_n|\le\frac{1}{n}$得$a_n\to0$。又$\sin1>0$，而$\pi<4<\frac{3\pi}{2}$，故$a_1>0$、$a_4<0$。应用定理\ref{cd:finite-extrema}，两种最值都能取得。直接证明时可取
\[
 \varepsilon=\dfrac12\min\{a_1,-a_4\}>0.
\]
当$n>N\ge4$时，$-\varepsilon<a_n<\varepsilon$，尾项严格小于$a_1$且大于$a_4$，故所有最值都可在前$N$项中找。对$b_n$，最大值为$b_1=1$；每项为正且趋0，故下确界为0，但没有最小值。\par
\textbf{方法点评} “收敛所以有界”只能保证上下确界存在；极限两侧的具体项才把确界提升为实际取得的最值。\par
\textbf{疑难点与变式} 结论只证明最值存在，没有声称$a_1$或$a_4$一定就是极值项。若只在极限一侧有项，另一种最值仍须另查。
\end{example}

\subsection{平移差与相等函数值}
\begin{method}{将相等函数值改为连续差的零点}
\textbf{适用条件} 要证明$f(t)=f(t+h)$，且两端点条件具有平移对称。\par
\textbf{步骤} 构造$G(t)=f(t)-f(t+h)$，在其定义区间端点计算符号；端点乘积非正时用介值定理。\par
\textbf{疑难点} 非正乘积允许零点在端点；只有严格异号才能保证零点在开区间。
\end{method}
\begin{example}{考研风格｜相隔定长的两个等高点}{难度：$\star\star$}\label{cd:ex:translate}
\textbf{题目} 设$a>0$，$f\in C[0,2a]$且$f(0)=f(2a)$，证明存在$t\in[0,a]$，使$f(t)=f(t+a)$。说明$t$能否总取在$(0,a)$。\par
\textbf{思路} 比较从左半区间和右半区间得到的同一差。\par
\textbf{详细解答} 令$G(t)=f(t)-f(t+a)$，则$G\in C[0,a]$，且
\[
 G(a)=f(a)-f(2a)=f(a)-f(0)=-G(0).
\]
若$G(0)=0$，取$t=0$；否则$G(0)G(a)=-G(0)^2<0$，介值定理给$t\in(0,a)$满足$G(t)=0$。开区间结论不总成立：取$f(x)=\sin(\frac{\pi x}{a})$，端点$f(0)=f(2a)=0$，但对$0<t<a$，
\[
 f(t)-f(t+a)=2\sin(\dfrac{\pi t}{a})>0,
\]
所以等高点只对应$t=0,a$。\par
\textbf{方法点评} 用差函数替代同时寻找两个点，可将几何等高问题变成一个未知数的零点问题。\par
\textbf{疑难点与变式} 若额外有$f(0)\ne f(a)$，端点严格异号，便能保证$t\in(0,a)$。
\end{example}

\subsection{归一化符号与无界区间零点}
\begin{method}{先在两端归一化，再截成闭区间}
\textbf{适用条件} $H$在$(0,\infty)$连续，能够分别在$0^+$和$+\infty$找到其符号。\par
\textbf{步骤} 用正的尺度因子消去趋零或增长的主量；由非零极限选出有限$x_1<x_2$使$H(x_1)H(x_2)<0$；最后在$[x_1,x_2]$使用介值定理。\par
\textbf{疑难点} 无界区间的两端不是可直接代入的函数值，必须实际选出有限端点。
\end{method}
\begin{example}{考研风格｜局部三阶信息强制出现零点}{难度：$\star\star\star$}\label{cd:ex:unbounded-root}
\textbf{题目} 设$f$在$(0,\infty)$连续有界，且
\[
 \lim_{x\to0^+}\dfrac{f(x)}{\sin x-x}=12.
\]
证明方程$f(x)+x^3=0$在$(0,\infty)$至少有一个实根。\par
\textbf{思路} 在0附近比较三阶主项，在无穷远用有界性压低$f$的阶。\par
\textbf{详细解答} 记$H(x)=f(x)+x^3$。由于$\frac{\sin x-x}{x^3}\to-\frac{1}{6}$，
\[
 \lim_{x\to0^+}\dfrac{H(x)}{x^3}
 =12\left(-\dfrac16\right)+1=-1.
\]
故存在$\delta>0$，使$0<x<\delta$时$H(x)<0$。另一方面，由$|f(x)|\le M$，
\[
 \left|\dfrac{f(x)}{x^3}\right|\le\dfrac{M}{x^3}\to0,
 \qquad \dfrac{H(x)}{x^3}\to1\quad(x\to+\infty).
\]
选$x_1\in(0,\delta)$，再取足够大的$x_2>x_1$使$H(x_2)>0$。$H$在$[x_1,x_2]$连续，介值定理给出所需零点。\par
\textbf{方法点评} 极限只负责找到符号，连续性负责连接这两个符号；两者的职责应分别写清。\par
\textbf{疑难点与变式} 本题没有单调性或导数信息，不能进一步断言根唯一。将12换成参数$c$时，上述论证在$c>6$仍成立。
\end{example}

\section{数列的单侧控制、方程列与隐式迭代}\label{cs:section}
\begin{examtip}
递推关系未必直接给出单调性。允许的小幅上升可用尾和补偿；方程定义的项需比较不同指标下的函数；当指数与指标一起变化时，应先给出统一界再取极限。
\end{examtip}

\subsection{用可和误差补偿单侧漂移}\label{cs:drift}
\begin{method}{把近似单调转成真正单调}
\textbf{条件} $x_n$有下界，$\varepsilon_n\ge0$，$\sum\varepsilon_n<\infty$，且$x_{n+1}\le x_n+\varepsilon_n$。
\par\textbf{步骤} 令$r_n=\sum_{k=n}^{\infty}\varepsilon_k$，研究$y_n=x_n+r_n$。由$r_n-r_{n+1}=\varepsilon_n$得$y_{n+1}\le y_n$；下界和$r_n\to0$给$x_n$收敛。
\par\textbf{疑难点} 仅有$\varepsilon_n\to0$不够，需要累计误差有限。
\end{method}
\begin{example}{考研风格｜不等式递推的补偿项}{难度：$\star\star$}\label{cs:ex:drift}
\textbf{题目} 正项数列满足
\[
x_{n+1}\le x_n+\dfrac{2^n}{(2^{n+1}-1)^2},\qquad n\ge1.
\]
证明$x_n$收敛，并给出统一上界。
\par\textbf{思路} 将误差放大成相邻倒数之差，主动构造补偿量。
\par\textbf{详细解答} 因$2^n-1<2^{n+1}-1$，
\[
\dfrac{2^n}{(2^{n+1}-1)^2}
\le\dfrac{2^n}{(2^n-1)(2^{n+1}-1)}
=\dfrac1{2^n-1}-\dfrac1{2^{n+1}-1}.
\]
于是$y_n=x_n+\frac{1}{2^n-1}$递减且为正，故$y_n\to L\ge0$，进而$x_n\to L$。
由$y_n\le y_1=x_1+1$，得到
\[
0<x_n\le x_1+1-\dfrac1{2^n-1}<x_1+1.
\]
\par\textbf{方法点评} 一般尾和不易写出时，可用更简单的$r_n\downarrow0$满足$r_n-r_{n+1}\ge\varepsilon_n$，证明完全相同。
\par\textbf{疑难点} 若仅要求$x_{n+1}-x_n\le\frac{1}{n}$，取$x_n=2+\sin(\ln n)$，由$\sin$的Lipschitz估计及$\ln(1+\frac{1}{n})\le\frac{1}{n}$可知条件成立，但沿$n_k=\lfloor\e^{2k\pi+\frac{\pi}{2}}\rfloor$与$m_k=\lfloor\e^{2k\pi+\frac{3\pi}{2}}\rfloor$分别趋3与1，数列不收敛。
\end{example}

\begin{example}{竞赛风格｜隐式不等式中的强制极限}{难度：$\star\star\star$}\label{cs:ex:forced}
\textbf{题目} $x_n>0$且$\ln x_n+\frac{1}{x_{n+1}}<1$。证明$x_n\to1$。
\par\textbf{思路} 从一个全局最小值构造相邻项的次序。
\par\textbf{详细解答} 设$g(t)=\ln t+\frac{1}{t}$，则$g'(t)=\frac{t-1}{t^2}$，故$g(t)\ge g(1)=1$，等号仅在$t=1$。
因此
\[
\ln x_n<1-\dfrac1{x_{n+1}}\le\ln x_{n+1},
\]
所以$x_n<x_{n+1}$。又$\frac{1}{x_{n+1}}>0$给$\ln x_n<1$，即$x_n<\e$。
数列递增有界，极限$L\ge x_1>0$。原不等式取极限得$\ln L+\frac{1}{L}\le1$，与$g(L)\ge1$合并，得$L=1$。
\par\textbf{方法点评} 上界可直接由原不等式读出；无需先假设发散再反证。
\par\textbf{疑难点} 正项极限原本可能为0；这里由递增性和$x_1>0$先排除0，才可代入$\ln L$与$\frac{1}{L}$。
\end{example}

\subsection{方程列中的统一余项与根的尺度}\label{cs:roots}
\begin{method}{先定位根，再处理随指标变化的幂}
\textbf{条件} 每个$n$有连续严格单调函数$F_n$与唯一根$x_n$。
\par\textbf{步骤} 用端点异号定位；计算$F_{n+1}(x_n)$判断根的移动；找到与$n$无关的$q<1$控制$x_n^n$，再求极限或误差。
\par\textbf{疑难点} 从$x_n\to1$不能推出$x_n^n\to1$；从每个$x_n<1$也不能直接推出$x_n^n\to0$。
\end{method}
\begin{example}{竞赛风格｜几何和方程的指数小误差}{难度：$\star\star\star$}\label{cs:ex:geometricroot}
\textbf{题目} 对$n\ge2$，设$x_n$为$\sum_{k=1}^{n}x^k=1$在$(0,1)$的根。证明$x_n\downarrow\frac{1}{2}$，并求
$\lim_{n\to\infty}2^n(x_n-\frac{1}{2})$。
\par\textbf{思路} 有限和先化闭式，根$x_2$为全部后续项提供统一界。
\par\textbf{详细解答} 多项式在$(0,1)$严格递增，取$x=\frac{1}{2}$时和为$1-2^{-n}<1$，取$x=1$时和为$n>1$，故唯一根在$(\frac{1}{2},1)$。
在$x_n$处增加一项后和大于1，所以$x_{n+1}<x_n$。特别地
\[
\dfrac12<x_n\le q:=x_2=\dfrac{\sqrt5-1}{2}<1.
\]
有限几何和给
\[
\dfrac{x_n(1-x_n^n)}{1-x_n}=1
\quad\Longleftrightarrow\quad 2x_n-1=x_n^{n+1}.
\]
右端不超过$q^{n+1}\to0$，故$x_n\to\frac{1}{2}$。
记$d_n=2x_n-1$，则$0<d_n\le q^{n+1}$，所以$nd_n\to0$，从而
$(2x_n)^n=\exp(n\ln(1+d_n))\to1$。于是
\[
2^n(x_n-\dfrac12)
=\dfrac{x_n}{2}(2x_n)^n\longrightarrow\dfrac14.
\]
\par\textbf{方法点评} 统一的$q<1$不只证明余项趋0，还给出$n$倍余项趋0，使指数尺度的计算合法。
\par\textbf{疑难点} 直接把“项数为$n$的和”中每一项换成极限，会把有限和变成未经证明的无穷和。
\end{example}
\begin{figure}[htbp]\centering
\begin{tikzpicture}
\begin{axis}[width=.88\textwidth,height=5.8cm,xmin=.45,xmax=.7,ymin=.65,ymax=1.4,
 axis lines=middle,xlabel={$x$},ylabel={$s_n(x)$},
 xtick={.5,.618},xticklabels={$\frac12$,$x_2$},ytick={1},
 legend style={font=\small,at={(.02,.98)},anchor=north west,draw=none},
 samples=81,domain=.45:.7]
\addplot[Teal,thick]{x+x^2};\addlegendentry{$n=2$}
\addplot[Gold,thick]{x+x^2+x^3};\addlegendentry{$n=3$}
\addplot[Ink,thick]{x+x^2+x^3+x^4+x^5};\addlegendentry{$n=5$}
\addplot[Rust,dashed]{1};\addlegendentry{$s_n=1$}
\draw[dashed,gray] (axis cs:.5,.65)--(axis cs:.5,1.4);
\end{axis}
\end{tikzpicture}
\caption{在固定$x\in(0,1)$处，$s_n(x)=\sum_{k=1}^nx^k$随$n$增加而增大，因此与水平线$y=1$的交点向左移动。所有$n\ge2$的交点都在$(\frac{1}{2},x_2]$内；这一几何夹持对应证明中的统一指数控制。}\label{cs:fig:rootfamily}
\end{figure}
\begin{method}{从图形识别根列的巧解入口}
横向看单个函数的单调性，确定根唯一；纵向看指标增加时整条曲线的移动，确定根的次序。几何直观提示结论后，应落实为$F_{n+1}(x_n)$的符号，而不依赖绘图分辨率。若统一界能把根限制在$q<1$以内，它还同时控制全部高次幂，这就是“定位”能够帮助“求尺度”的原因。
\end{method}

\begin{example}{竞赛拓展｜接近1的根与对数尺度}{难度：$\star\star\star$}\label{cs:ex:logroot}
\textbf{题目} $x_n\in(0,1)$满足$\ln x_n+x_n^n=0$。证明根唯一，并求
$\lim_{n\to\infty}\frac{n(1-x_n)}{\ln n}$。
\par\textbf{思路} 将未知幂转换为指数方程，先估计$-\ln x_n$。
\par\textbf{详细解答} $\ln x+x^n$在$(0,1)$的导数为$\frac{1}{x}+nx^{n-1}>0$，在左端趋$-\infty$，在1处为1，故根唯一。
设$y_n=-\ln x_n>0$，则$y_n=\e^{-ny_n}$。记$t_n=ny_n$，得到
\[
t_n\e^{t_n}=n,\qquad t_n+\ln t_n=\ln n.
\]
当$n$充分大时$t_n>1$，于是
\[
\ln n-\ln\ln n\le t_n\le\ln n.
\]
故$\frac{t_n}{\ln n}\to1$，$y_n=\frac{t_n}{n}\to0$，且
\[
\dfrac{n(1-x_n)}{\ln n}
=\dfrac{1-\e^{-y_n}}{y_n}\dfrac{t_n}{\ln n}\longrightarrow1.
\]
若用Lambert函数表示，精确式为$x_n=\exp[-\frac{W(n)}{n}]$；上面的夹逼已经独立证明所需渐近关系。
\par\textbf{方法点评} 变量$n(1-x_n)$本身发散，对数归一化才能给出有限非零极限。
\par\textbf{疑难点} 不能将$x_n^n$写成$(\lim x_n)^n$；应在$y_n\e^{ny_n}=1$中保留$n$。
\end{example}

\begin{example}{考研风格｜负根的一阶边界层}{难度：$\star\star\star$}\label{cs:ex:negativeroot}
\textbf{题目} $a>0$，$x_n$为$\e^{ax}+x^{2n+1}=0$在$(-1,0)$的根。求$\lim x_n$及$\lim n(x_n+1)$。
\par\textbf{思路} 单调性给唯一根，对正量$-x_n$取对数。
\par\textbf{详细解答} 函数在$-1$处值为$\e^{-a}-1<0$，在0处为1，导数$a\e^{ax}+(2n+1)x^{2n}>0$，故根唯一。
令$y_n=-x_n\in(0,1)$，方程等价于
\[
(2n+1)\ln y_n=-ay_n.
\]
右边有界，故$\ln y_n\to0$，$y_n\to1$，即$x_n\to-1$。再由
$-\ln y_n\sim1-y_n$，得到
\[
n(x_n+1)=n(1-y_n)
 =\dfrac{1-y_n}{-\ln y_n}\dfrac{nay_n}{2n+1}
 \longrightarrow\dfrac a2.
\]
\par\textbf{方法点评} 有界性直接控制归一化对数，求根极限不必先证明整个根列单调。
\par\textbf{疑难点} 对$x_n$本身取实对数无意义，须先转换为正量$y_n$。
\end{example}

\subsection{反向递推与正确的几何缩放}\label{cs:inverse}
\begin{method}{先证明每一步可解，再寻找渐近乘子}
\textbf{条件} 递推写作$x_n=H(x_{n+1})$，$H$在所用区间连续严格递增。
\par\textbf{步骤} 确定$H$的值域，确定唯一后继；比较$H(t)$与$t$；若$x_n\to0$且$\frac{H(t)}{t}\to\lambda>1$，则$\frac{x_{n+1}}{x_n}\to\frac{1}{\lambda}$。
\par\textbf{疑难点} 比值极限只给几何主尺度；证明$\lambda^nx_n$有有限非零极限还需控制比值误差的累加。
\end{method}
\begin{example}{竞赛风格｜反向对数递推的缩放常数}{难度：$\star\star\star$}\label{cs:ex:inverselog}
\textbf{题目} $x_1=1$，对$n\ge1$由
$x_n=x_{n+1}+2\ln(1+x_{n+1})$定义正项数列。证明$x_n\to0$、$2^nx_n\to0$，并证明$3^nx_n$趋于某个有限正数。
\par\textbf{思路} 先求线性化比例，再用对数乘积证明精确缩放的收敛。
\par\textbf{详细解答} $H(t)=t+2\ln(1+t)$在$[0,\infty)$严格递增，$H(0)=0$且趋于无穷，每步有唯一正解，满足$0<x_{n+1}<x_n$。
极限$L\ge0$满足$2\ln(1+L)=0$，故$L=0$。由
\[
\dfrac{x_{n+1}}{x_n}
=\dfrac{t}{t+2\ln(1+t)}\bigg|_{t=x_{n+1}}\longrightarrow\dfrac13,
\]
可选$q\in(\frac{1}{3},\frac{1}{2})$，使充分大的$n$满足$x_{n+1}\le qx_n$。因此$\sum x_n<\infty$，且$2^nx_n\to0$。
记$z_n=3^nx_n$。由$\ln(1+t)<t$得$\frac{3t}{H(t)}>1$，故$z_n$递增。又Taylor展开给
\[
H(t)=3t-t^2+O(t^3),\qquad
\ln\dfrac{z_{n+1}}{z_n}=\ln\dfrac{3t}{H(t)}
=\dfrac t3+O(t^2),\quad t=x_{n+1}.
\]
因$\sum x_n<\infty$，对数增量之和收敛，故$\ln z_n$趋于有限实数，$z_n\to C\in(0,\infty)$。
\par\textbf{方法点评} 先用比值法取得粗几何上界，再将其用于更精确的乘积收敛，形成两步估计。
\par\textbf{疑难点} 不能由$\frac{z_{n+1}}{z_n}\to1$推出$z_n$收敛；$\ln z_n$的可和增量才是关键。
\end{example}

\begin{example}{竞赛拓展｜高次开方的全局Newton迭代}{难度：$\star\star\star$}\label{cs:ex:newtonpower}
\textbf{题目} $m\ge2$为整数，$A>0$，$x_1>0$，令
$x_{n+1}=\frac{(m-1)x_n+\frac{A}{x_n^{m-1}}}{m}$。证明$x_n\to r=A^{\frac{1}{m}}$，并求非恒定情形的二次误差系数。
\par\textbf{思路} 将$m-1$个$x_n$与一个$\frac{A}{x_n^{m-1}}$作算术几何平均。
\par\textbf{详细解答} 均值不等式给$x_{n+1}\ge r$。从$n\ge2$起，
\[
x_{n+1}-x_n=\dfrac{A-x_n^m}{mx_n^{m-1}}\le0.
\]
故从第二项起递减有下界，极限$L\ge r>0$，代入递推得$L^m=A$，即$L=r$。
设$T(x)=\frac{(m-1)x+Ax^{1-m}}{m}$，计算
\[
T'(r)=0,\qquad T''(r)=\dfrac{m-1}{r}.
\]
Taylor公式给
$x_{n+1}-r=\frac{m-1}{2r}(x_n-r)^2+o((x_n-r)^2)$。
若$x_1\ne r$，均值不等式严格，全部误差在适当指标后为正，故
\[
\lim_{n\to\infty}\dfrac{x_{n+1}-r}{(x_n-r)^2}=\dfrac{m-1}{2r}.
\]
\par\textbf{方法点评} 同一递推同时具有全局收敛保证与局部二次速度。
\par\textbf{疑难点} 当$x_1<r$时首步上升，不能声称从第一项起就递减；$x_1=r$时误差比是$\frac{0}{0}$，应单独说明为恒定解。
\end{example}
\begin{warning}
数列题的三个层次应分开：递推是否定义了实数后继、极限是否存在、所选缩放是否给出有限非零极限。先解不动点方程或仅算比值，通常只能处理其中一部分。
\end{warning}

\section*{章末分层习题}\addcontentsline{toc}{section}{章末分层习题}
\refstepcounter{exercise}\label{q:1:1}\par\noindent\textbf{习题\theexercise\quad 基础题} 求$\lim_{x\to0}\frac{\sin5x}{x}$。\quad\hyperref[a:1:1]{[答案]}\par
\refstepcounter{exercise}\label{q:1:2}\par\noindent\textbf{习题\theexercise\quad 基础题} 判断$\frac{|x|}{x}$在0的极限。\quad\hyperref[a:1:2]{[答案]}\par
\refstepcounter{exercise}\label{q:1:3}\par\noindent\textbf{习题\theexercise\quad 基础题} 求$\lim_{n\to\infty}\frac{3n+1}{2n-1}$。\quad\hyperref[a:1:3]{[答案]}\par
\refstepcounter{exercise}\label{q:1:4}\par\noindent\textbf{习题\theexercise\quad 考研风格} 求$\lim_{x\to0}\frac{\cos x-1+\frac{x^2}{2}}{x^4}$。\quad\hyperref[a:1:4]{[答案]}\par
\refstepcounter{exercise}\label{q:1:5}\par\noindent\textbf{习题\theexercise\quad 考研风格} 求$\lim_{n\to\infty}\sum_{k=1}^n \frac{k}{n^2}$。\quad\hyperref[a:1:5]{[答案]}\par
\refstepcounter{exercise}\label{q:1:6}\par\noindent\textbf{习题\theexercise\quad 考研风格} 给$\frac{\sin x}{x}$在0补值使其连续。\quad\hyperref[a:1:6]{[答案]}\par
\refstepcounter{exercise}\label{q:1:7}\par\noindent\textbf{习题\theexercise\quad 竞赛拓展} 求$\lim_{n\to\infty}n\ln(1+\frac{1}{n})$。\quad\hyperref[a:1:7]{[答案]}\par
\refstepcounter{exercise}\label{q:1:8}\par\noindent\textbf{习题\theexercise\quad 竞赛拓展} 设$a_1=1,a_{n+1}=\frac{a_n+\frac{2}{a_n}}{2}$，证明收敛并求极限。\quad\hyperref[a:1:8]{[答案]}\par
\refstepcounter{exercise}\label{q:1:9}\par\noindent\textbf{习题\theexercise\quad 竞赛拓展} 求$\lim_{n\to\infty}\frac1n\sum_{k=1}^n\sqrt{\frac{k}{n}}$。\quad\hyperref[a:1:9]{[答案]}\par
\section*{章末方法速查}\addcontentsline{toc}{section}{章末方法速查}
\par\noindent\begin{tabularx}{\textwidth}{p{23mm}XXX}\toprule 方法&适用条件&典型信号&易错点\\\midrule
等价替换&乘积与商&一阶主项&相消丢阶\\Taylor&局部光滑&多项相消&余项未除分母\\Stolz&分母递增无界&求和比&遗漏假设\\积分和&连续或另有估计&步长与采样点&奇点未经处理\\
\bottomrule\end{tabularx}
